

\subsection{From extended affine to affine.}\label{subsec extended affine to affine}
Most, if not all, of the content in this subsection is well-known to experts, but this subsection is warranted by the lack of a consolidated reference in the literature (as far as the authors are aware).
We retain the notation as in \cref{subsec extended affine Weyl groups}:
$W$ denotes the affine Weyl group, $W_{\ext}=\Omega \ltimes W $ denotes the extended affine Weyl group of $\PGL_n$, $\omega$ is a generator of $\Om \subset W_{\ext}$, with $\omega s_i \omega^{-1}=s_{i+1}$, and $\omega^n=1$, $\varpi_k$ denotes the fundamental weight for $1\leq k \leq n-1$, $\rot: \Lambda_{wt} \rightarrow \Om$ denotes the quotient map with respect to the root lattice.
While $W_{\ext}$ is not a Coxeter group, the length function on $W$ extends to a length function on $W_{\ext}$ such that $\ell(\omega^k)=0$ for any $k$.

% Explicitly, the translation $t_{\omega_k}$ by $\omega_k$ and $\varpi_k$ are related by
% $$t_{\omega_k}=\varpi_k w_{\hh{0}}w_{\hh{0,k}}.$$
% \PT{Cite MacDonald, Affine Hecke Algebras and Orthogonal Polynomials book, section 2.5? Or do you have a preferred reference?}


Since $W_{\ext}= W_{\hh{0}}\ltimes\Lambda_{\wt}$, double cosets in
$W_{\hh{0}}\backslash W_{\ext}/W_{\hh{0}}$ have a unique representative in $\Lambda_{\wt}/W_{\hh{0}}$ which is naturally in bijection with $\Lambda_{\wt}^+$ . 
We refer to elements of $W_{\hh{0}}\backslash W_{\ext}/W_{\hh{0}}$  as \emph{spherical double cosets}, and the double coset $W_{\hh{0}}t_{\lambda}W_{\hh{0}}$ as the \emph{spherical double coset of $\lambda$}.
Note that when $\lambda \in \Lambda_{\rt}^+$, the spherical double  coset of $\lambda$ is an element of $W_{\hh{0}}\backslash W/W_{\hh{0}} \subset W_{\hh{0}}\backslash W_{\ext}/W_{\hh{0}}$.

\begin{defn}\label{defn the map Xi}
    We define the map 
    \begin{equation}\Xi: \Lambda_{\wt}^+ \rightarrow \coprod_{0 \leq k\leq n-1} W_{\hh{0}}\backslash W/W_{\hh{k}}\end{equation}
    by setting 
    \begin{equation}\Xi(\lambda):=W_{\hh{0}}t_{\lambda}\omega^{-\rot(\lambda)}W_{\hh{\rot(\lambda)}}
    \end{equation}   
    so that 
    \begin{equation}
        W_{\hh{0}}t_{\lambda}W_{\hh{0}}= \Xi(\lambda)\omega^{\rot(\lambda)}.
    \end{equation}
\end{defn}

\begin{lemma}\label{lemma Xi vs psi}
The map $\Xi$ is a bijective correspondence
% \begin{align}
%    \Lambda_{\wt}^+ \longleftrightarrow W_{\hh{0}}\backslash &W_{\ext}/W_{\hh{0}} \longleftrightarrow  \coprod_{k \in \Om}W_{\hh{0}}\backslash W/W_{\hh{k}}\\
%    \notag \lambda \longleftrightarrow W_{\hh{0}}&t_\lambda W_{\hh{0}} \longleftrightarrow W_{\hh{0}}t_{\lambda}\omega^{-\rot{\lambda}}W_{\hh{k}}
% \end{align}
which relates to the correspondence $\psi_0$ in \cref{def:psia} by the equation
$$\inv\circ\psi_0(\lambda)=\Xi(\lambda),$$
where $\inv: \displaystyle \coprod_{k, l \in \Om} W_{\hh{l}}\backslash W/W_{\hh{k}} \xrightarrow{\sim} \coprod_{l,m \in \Om}W_{\hh{k}}\backslash W/W_{\hh{l}}$ is the isomorphism induced by the map $w \mapsto w^{-1}$.
\end{lemma}
\begin{proof}
    This is an immediate consequence of \cref{lemma coset correspondence for extended affine spherical cosets}.
\end{proof}

\begin{remark}\label{rmk coset mismatch}
    For $\lambda \in \Lambda_{\wt}^+$ with $\rot(\lambda)=k$, the spherical double cosets $\omega^{-k}\psi_0(\lambda)$ and $\Xi(\lambda)\omega^k$ coincide with $W_{\hh{0}}t_{-\lambda}W_{\hh{0}}$ and $W_{\hh{0}}t_{\lambda}W_{\hh{0}}$ respectively. 
    Since the dominant weight representative in the $W$-orbit of $-\lambda$ is $-w_{\hh{0}}(\lambda)$,  the dominant representative of the spherical double coset $\omega^{-k}\psi_0(\lambda)$ is $-w_0(\lambda)$.
\end{remark}
% \begin{example}\label{example coset for fundamental weights}
%     Under the correspondence above, the $(W_{\hh{0}},W_{\hh{0}})$-double coset $ W_{\hh{0}}t_{\omega_k}W_{\hh{0}}$ corresponds to the $(W_{\hh{0}},W_{\hh{k}})$-double coset $ W_{\hh{0}}1W_{\hh{k}}$.
% \end{example}

Let $\HH_{\ext}$ denote the extended affine Hecke algebra associated to $W_{\ext}$, which has the affine Hecke algebra $\HH_{\aff}$ as a subalgebra. 
To make the formulation of the Satake isomorphism in terms of the Hecke algebroid precise, we compare the presentations of the Hecke algebra in \cite{WillSingular}  and \cite{Ach} (which offers a consolidated alternative to \cite{LusztigGS}).

The presentation we follow in \cref{subsec-Hecke} is from \cite{WillSingular}, so we start by recalling the explicit conventions from the same. 
The Hecke algebra in \cite{WillSingular} is defined using the quadratic relation
\begin{equation}H_s^2=(v^{-1}-v)H_s + 1.
\end{equation}
For a double coset $p \in W_I\backslash W/W_J$ with maximal element $\bar{p}$, the corresponding explicit standard basis element and Kazhdan-Lusztig basis element are given by
\begin{equation}\label{eq Williamson std basis and KL basis Hecke algebroid}
 \HAB{I}{J}{p}:=  \sum_{x\in p}v^{\ell(\bar{p})-\ell(x)}H_x, \qquad
 \KLB{I}{J}{p}:=\KLB{}{}{\bar{p}}.
\end{equation}
Note that the standard basis here is normalized precisely so that the transformation matrix from the standard basis to the Kazhdan-Lusztig basis (of $\Hecke{I}{J}$) is unipotent.

There is an action of $\Omega$ on $\HH_{\aff}$, corresponding to the action of $\Omega$ on the affine Dykin diagram, and $\HH_{\ext}$ may be thought of as the twisted group ring
\begin{equation}\label{eq H-ext as a semidirect product of H-aff and omega}
\HH_{\ext}=\HH_{\aff} \rtimes \Omega.
\end{equation}
Indeed, since $\ell(\omega^k)=0$, the elements $H_{\omega^k}\in \HH_{\ext}$ satisfies the equations 
\begin{equation}\label{eq Hecke multiplication length zero elements}
  H_x H_{\omega^k}=H_{x\omega^k}, \quad H_{\omega^k} H_x=H_{\omega^k x}  , \quad H_{\omega^k}H_x=H_{\omega^k(x)}H_{\omega^k}
\end{equation} for any $x\in W$ where $\omega^k(x)=\omega^k x \omega^{-k}$.

Now we recall the presentation for the extended affine Hecke algebra and explicit formulas for the standard basis for the corresponding spherical Hecke algebra from \cite{Ach}.
This presentation for the extended affine Hecke algebra, as in \cite{Ach}, is given over $\Z[q^{\pm1/2}]$, with basis elements $T_w$ for $w \in W_{\ext}$, and the quadratic relation 
\begin{equation}\label{eq Hecke algebra presentation Williamson vs Lusztig}
    T_s^2=(q-1)T_s+q T_1.
\end{equation}
Comparing with the presentation  from \cite{WillSingular}, we have
\begin{equation}v = q^{-1/2}, \qquad H_w=v^{\ell(w)}T_w.
\end{equation}



The spherical Hecke algebra $\HH_{\sph}$ is then defined as a subalgebra of the extended affine Hecke algebra localized at the Coxeter polynomial 

\begin{equation}\label{eq coxeter polynomial achar to williamson}
W_q:=\sum_{w \in W_{\hh{0}}}q^{\ell(w)}=\sum_{w \in W_{\hh{0}}}v^{-\ell(w)}=\tilde{\pi}(\hh{0})=v^{-\ell(w_{\hh{0}})}\pi(\hh{0})
\end{equation}
(for $\tilde{\pi}, \pi$ is as in \cref{subsec-Hecke}).
Explicitly, this subalgebra is the $\Z[v^{\pm 1}]$-span of $\set{K_\lambda}_{\lambda \in \lambda_{\wt}^+}$, where
\begin{equation}\label{eq standard basis spherical Hecke}
    K_\lambda:=W_q^{-1} \sum_{w\in W_{\hh{0}}t_\lambda W_{\hh{0}}}T_w.
\end{equation}
This subalgebra also coincides with $K_0\HH_{\ext}\cap \HH_{\ext}K_0=K_0\HH_{\ext}K_0$ (\cite[Lemma 9.3.6]{Ach}).
Note here that $K_0$ is an idempotent, and is explicitly related to the quasi-idempotent $\un{H}_{w_{\hh{0}}}$ as
\begin{equation}
    K_0=q^{\ell(w_{\hh{0}})/2}W_{q}^{-1}\un{H}_{w_{\hh{0}}}=\frac{1}{\pi(\hh{0})}\un{H}_{w_{\hh{0}}}.
\end{equation}
The spherical Hecke algebra $\HH_{\sph}$ inherits the bar involution from $\HH_{\ext}$, and there is an invariant basis $\set{\un{S}_{\lambda}}_{\lambda \in \Lambda^+_{\wt}}$ unique up to certain properties, as in \cite[Theorem 9.3.7]{Ach}.


For $\lambda \in \Lambda_{\wt}^+$, let $L_{\lambda}$ denote the irreducible representation of highest weight $\lambda$ in $\Rep_{\mathbb{C}}(\mathfrak{sl}_n)$.
Lusztig, in \cite{LusztigGS}, showed that the Satake isomorphism precisely identifies $\ch(L_\lambda)$ with $\un{S}_\lambda$. We formally state these results below (see \cite[Theorem 9.4.1, 9.4.2]{Ach}).

\begin{theorem}[Satake]\label{thm Satake original}
    There is an isomorphism of $\Z[v^{\pm 1}]$-algebras 
    $$\gamma: \HH_{\sph} \xrightarrow{\sim} \Z[v^{\pm 1}][\Lambda_{\wt}]^W.$$
\end{theorem}

\begin{theorem}[Lusztig]\label{thm Lusztig Satake}
    The Satake isomorphism $\gamma$ identifies $\un{S}_\lambda$ with $\ch(L_\lambda)$.
\end{theorem}


Inside $\HH_{\sph}$, we have the subalgebra $\HH_{\sph, \aff}$ spanned by the $K_\lambda$ for $\lambda \in \Lambda_{\rt}^+$.
The decomposition in \cref{eq H-ext as a semidirect product of H-aff and omega} then induces a decomposition
\begin{equation}
    \HH_{\sph}=\HH_{\sph, \aff} \rtimes \Om.
\end{equation}

The following result is well known (cf. \cite[Remark 2.2.2 (3)]{WillThesis}), but we provide a proof for the sake of clarity.

\begin{lemma}\label{lemma affine spherical Hecke in Hecke algebroid}
 There is an isomorphism of $\Z[v^{\pm 1}]$-algebras   
 $$\HH_{\sph, \aff} \cong \Hecke{\hh{0}}{\hh{0}}$$
 which sends the basis elements $K_\lambda$ to $v^{-\langle2\rho^\vee, \lambda \rangle}\HAB{\hh{0}}{\hh{0}}{W_{\hh{0}}t_{\lambda}W_{\hh{0}}}$ and $\un{S}_\lambda$ to $\KLB{\hh{0}}{\hh{0}}{W_{\hh{0}}t_{\lambda}W_{\hh{0}}}$.
\end{lemma}
\begin{proof}
For any $\lambda \in \Lambda_{\wt}$, the maximal coset representative of $W_{\hh{0}}t_\lambda W_{\hh{0}}$ has length $\ell(w_{\hh{0}}) + \langle 2\rho^\vee, \lambda \rangle$ (see \cite[Chapter 2]{Macdonaldhecke}\footnote{This fact may also be seen as a consequence of \cref{lemma dist coset reps}, \cref{thm:weightsadd} and \cref{lemma distinguished reps for fundamental weights}.}).
From \cref{eq standard basis spherical Hecke}, \cref{eq coxeter polynomial achar to williamson} and \cref{eq Hecke algebra presentation Williamson vs Lusztig},
we have that for $\lambda \in \Lambda_{\rt}^+$,
$$K_\lambda= \frac{v^{\ell(w_{\hh{0}})}}{\pi(\hh{0})} \sum_{w\in W_{\hh{0}}t_\lambda W_{\hh{0}}}v^{-\ell(w)}H_w= \frac{v^{-\langle2 \rho^\vee, \lambda \rangle}}{\pi(\hh{0})}\; \HAB{\hh{0}}{\hh{0}}{W_{\hh{0}}t_\lambda W_{\hh{0}}}$$
in $\HH_{\aff}[W_q^{-1}]=\HH_{\aff}[\pi(\hh{0})^{-1}].$
Additionally, given $\mu \in \Lambda_{\rt}$, we have
$$K_\lambda K_\mu
=\frac{v^{-\langle 2 \rho^\vee, \lambda+\mu \rangle}}{\pi(\hh{0})^2} \HAB{\hh{0}}{\hh{0}}{W_{\hh{0}}t_\lambda W_{\hh{0}}} \HAB{\hh{0}}{\hh{0}}{W_{\hh{0}}t_\mu W_{\hh{0}}}
=\frac{1}{\pi(\hh{0})}(v^{-\langle 2\rho^\vee, \lambda \rangle} \HAB{\hh{0}}{\hh{0}}{\lambda})
\ast_{\hh{0}}(v^{-\langle 2\rho^\vee, \mu \rangle} \HAB{\hh{0}}{\hh{0}}{\mu}).$$
It follows that the $\Z[v^{\pm 1}]$-linear automorphism of $\HH_{\aff}[\pi(\hh{0})^{-1}]$ defined by multiplication by $\pi(\hh{0})$ identifies  $K_\lambda$ with $v^{-\langle 2\rho^\vee, \lambda \rangle} \HAB{\hh{0}}{\hh{0}}{\lambda}$, and induces an isomorphism of algebras $\HH_{\sph, \aff} \cong \Hecke{\hh{0}}{\hh{0}}$.
This isomorphism preserves the bar involution, and the statement for the invariant basis may be deduced using the defining properties, or by using \cite[Theorem 6.12]{LusztigGS}, which states that
$\un{S}_\lambda$ (denoted by $C'_\lambda$ in \cite{LusztigGS}) equals 
$$\frac{1}{\pi(\hh{0})}\KLB{\hh{0}}{\hh{0}}{W_{\hh{0}}t_\lambda W_{\hh{0}}}$$
in $\HH_{\ext}[\pi(\hh{0})^{-1}]$.
\end{proof}

Now consider an algebra $\HHHH$ made using certain $\Hom$-spaces in the Hecke algebroid $\HH$ as defined below. We follow the notation from \cref{subsec-Hecke}.
\begin{defn}
  Let  $\HHHH:= \bigoplus_{0 \leq k \leq n-1}\Hecke{\hh{0}}{\hh{k}}$. 
  Define an $\Om-$graded $\Z[v,v^{-1}]$-algebra structure 
  $\bullet:\HHHH \times \HHHH \rightarrow \HHHH$ such that
  for $f \in \Hecke{\hh{0}}{\hh{k}}$, $g \in \Hecke{\hh{0}}{\hh{l}}$,
\begin{equation} \label{eq multiplication in HHHH}
    f\bullet g:= f\ast_{\hh{k}}\omega^k(g) \in \Hecke{\hh{0}}{\hh{k+l}},
\end{equation}
extended linearly to $\HHHH$. 
\end{defn}

\begin{lemma}\label{lemma HHHH isomorphic to spherical Hecke algebra}
  There is an isomorphism of $\Om$-graded $\Z[v^{\pm 1}]$-algebras
  $\HH_{\sph} \cong \HHHH$
  which sends the basis elements $K_\lambda$ to $v^{-\langle2 \rho^\vee, \lambda \rangle} \HAB{\hh{0}}{\hh{k}}{\Xi(\lambda)}$ and $\un{S}_\lambda$ to $\KLB{\hh{0}}{\hh{k}}{\Xi(\lambda)},$ for $\lambda \in \Lambda_{\wt}^+$ with $\rot(\lambda)=k$. In particular, $\HHHH$ is a commutative algebra.
\end{lemma}
\begin{proof}
    We proceed similarly as in the proof of \cref{lemma affine spherical Hecke in Hecke algebroid}, and use the equations in \cref{eq Hecke multiplication length zero elements} to see that  
    given $\lambda \in \Lambda_{\wt}^+$ with $\rot(\lambda)=k$, we have
    $$K_\lambda
    = \frac{v^{\ell(w_{\hh{0}})}}{\pi(\hh{0})} \sum_{w\in W_{\hh{0}}t_\lambda W_{\hh{0}}}v^{-\ell(w)}H_w 
    = \frac{v^{\ell(w_{\hh{0}})}}{\pi(\hh{0})} \sum_{w\in \Xi(\lambda)}v^{-\ell(w)}H_w H_{\omega^k}
    = \frac{v^{-\langle2 \rho^\vee, \lambda \rangle}}{\pi(\hh{0})} \HAB{\hh{0}}{\hh{k}}{\Xi(\lambda)} H_{\omega^k}$$
    in $\HH_{\ext}[\pi(\hh{0})^{-1}]$.
    Additionally, given $\mu \in \Lambda_{\wt}^+$ with $\rot(\mu)=l$, we have
    \begin{align*}
    K_\lambda K_\mu
    =\frac{v^{-\langle 2 \rho^\vee, \lambda+\mu \rangle}}{\pi(\hh{0})^2} \HAB{\hh{0}}{\hh{k}}{\Xi(\lambda)}H_{\omega^k} &\HAB{\hh{0}}{\hh{l}}{\Xi(\mu)}H_{\omega^l}
    =\frac{v^{-\langle 2 \rho^\vee, \lambda+\mu \rangle}}{\pi(\hh{0})\pi(\hh{k})} \HAB{\hh{0}}{\hh{k}}{\Xi(\lambda)}\HAB{\hh{k}}{\hh{k+l}}{\omega^k({\Xi(\mu)})}H_{\omega^{k+l}}\\
    &=\frac{1}{\pi(\hh{0})}(v^{-\langle 2\rho^\vee, \lambda \rangle} \HAB{\hh{0}}{\hh{k}}{\Xi(\lambda)})
    \ast_{\hh{k}}\omega^k \left(v^{-\langle 2\rho^\vee, \mu \rangle} \HAB{\hh{0}}{\hh{l}}{\Xi(\mu)}\right)H_{\omega^{k+l}}\\
    &=\frac{1}{\pi(\hh{0})}\left((v^{-\langle 2\rho^\vee, \lambda \rangle} \HAB{\hh{0}}{\hh{k}}{\Xi(\lambda)})
    \bullet (v^{-\langle 2\rho^\vee, \mu \rangle} \HAB{\hh{0}}{\hh{l}}{\Xi(\mu)})\right)H_{\omega^{k+l}}.
\end{align*}
It follows that the $\Z[v^{\pm 1}]$-linear isomorphism sending $K_\lambda$ to $v^{-\langle2 \rho^\vee, \lambda \rangle} \HAB{\hh{0}}{\hh{k}}{\Xi(\lambda)}$ is an isomorphism of $\Om$-graded algebras.
The statement for $\un{S}_{\lambda}$ then follows from \cite[Theorem 6.12]{LusztigGS} which implies that for $\lambda \in \Lambda_{\wt}^+$,
$$\un{S}_{\lambda}= \frac{1}{\pi(\hh{0})}\KLB{\hh{0}}{\hh{k}}{\Xi(\lambda)}H_{\omega^k}$$
in $\HH_{\ext}[\pi(\hh{0})^{-1}]$.
The statement regarding commutativity follows from the commutativity of $\HH_{\sph}.$
\end{proof}

To match with the conventions in this paper (in particular, the conventions in the coset correspondence in \cref{def:psia} and reading diagrams from right to left), we use the anti-involution of the Hecke algebroid, induced by that of the Hecke algebra which sends $H_w$ to $H_{w^{-1}}$.
Note that this lifts the involution of $W$ given by $w \mapsto w^{-1}$,  and identifies $\HAB{\hh{k}}{\hh{l}}{p}$ with $\HAB{\hh{l}}{\hh{k}}{\iota(p)}$ (where iota is as in \cref{lemma Xi vs psi}).
This anti-involution further identifies the commutative algebra $\HHHH^{\op}=\HHHH$ with $\HHH$ (see \cref{defn the algebra HHH}), identifying $\HAB{\hh{0}}{\hh{k}}{\Xi(\lambda)}$ with $\HAB{\hh{k}}{\hh{0}}{\Psi_0(\lambda)}$.
Combining \cref{thm Satake original}, \cref{thm Lusztig Satake} and \cref{lemma HHHH isomorphic to spherical Hecke algebra} with this observation, we then have the following theorem.
\begin{theorem}\label{thm Satake isomorphism to HHH}
 There is an isomorphism of $\Om$-graded $\Z[v^{\pm 1}]$-algebras
 \begin{gather}
     \notag \Z[v^{\pm 1}][\Lambda_{\wt}]^W \xrightarrow{\sim} \HH_{\sph} \xrightarrow{\sim} \HHHH \xrightarrow{\sim}\HHH \\
     \notag  \ch(L_\lambda) \mapsto \un{S}_\lambda \mapsto \KLB{\hh{0}}{\hh{\rot(\lambda)}}{\Xi(\lambda)} \mapsto \KLB{\hh{\rot(\lambda)}}{\hh{0}}{\Psi_0(\lambda)}. 
 \end{gather}   
\end{theorem}



% Combining the results from \cite{LusztigGS}, we then have:
% \begin{lemma}
%     The maps $\ast_{\hh{a+k}}: \Hecke{\hh{a}}{\hh{a+k}}\times \Hecke{\hh{a+k}}{\hh{a+k+l}} \rightarrow \Hecke{\hh{a}}{\hh{a+k+l}}$ can be read off from the multiplication rules within the representation ring of $\mathbf{sl}_n$.
%     More precisely, 
%     $$\ch(L_\lambda \otimes L_\mu)\overset{\text{Satake}}{=}S_\lambda S_\mu=\pi(\hh{0})^{-1}\;\KLB{\hh{a}}{\hh{a+k}}{\varpi_a(\tilde{p}_\lambda)} \ast_{\hh{a+k}} \KLB{\hh{a+k}}{\hh{a+k+l}}{\varpi_{a+k}(\tilde{p}_\mu)}H_{\varpi_{a+k+l}}.$$
% \end{lemma}
% \begin{example}\label{example hecke conv square flop dimension computation}
%     In $\Rep(\mathfrak{sl}_n)$, we have that $L_{\omega_k} \otimes L_{\omega_1}= L_{\omega_{k+1}} \oplus L_{\omega_k+\omega_1}$. 
%     The above lemma then tells us that $\KLB{\hh{a}}{\hh{a+k}}{\varpi_a(\tilde{p}_{\omega_k})} \ast_{\hh{a+k}} \KLB{\hh{a+k}}{\hh{a+k+1}}{\varpi_{a+k}(\tilde{p}_{\omega_1})}=\KLB{\hh{a}}{\hh{a+k}}{\varpi_a(\tilde{p}_{\omega_{k+1}})}+ \KLB{\hh{a+k}}{\hh{a+k+1}}{\varpi_{a+k}(\tilde{p}_{\omega_k+\omega_1})}$.
%     In particular, when $a=0,$ we have
%     $$\KLB{\hh{0}}{\hh{k}}{\tilde{p}_{\omega_k}} \ast_{\hh{k}} \KLB{\hh{k}}{\hh{k+1}}{\varpi_{k}(\tilde{p}_{\omega_1})}=\KLB{\hh{0}}{\hh{k}}{\tilde{p}_{\omega_{k+1}}}+ \KLB{\hh{k}}{\hh{k+1}}{\varpi_{k}(\tilde{p}_{\omega_k+\omega_1})}.$$
%     Note that from \cref{example coset for fundamental weights}, we know that the double cosets corresponding to the fundamental weights in $W_{\aff}$ are the cosets containing 1.
% \end{example}


\PTv2{Leaving the following comment for v2 modifications...}
\PTv2{I think I now understand the reason for the mismatch. The Soergel version of geometric Satake we construct is actually a composition of usual geometric Satake followed by the inversion anti-automorphism on the loop group/affine Kac Moody group. Normal geometric Satake at the affine Kac moody group level would take $$L_\lambda \mapsto IC_\lambda \mapsto IC_{W_{\hat{0}}1W_{\hat{k}}} \in \Perv_{P_{\hat{0}}\times P_{\hat{\hat{k}}}}(G_{KM}).$$
Composing with the isomorphism $$\Perv_{P_{\hat{0}}\times P_{\hat{\hat{k}}}}(G_{KM}) \xrightarrow{\sim} \Perv_{P_{\hat{k}}\times P_{\hat{0}}}(G_{KM})$$
induced by the inverse map of $G_{KM}$, we get our version. 
At the bimodule level, this is likely swapping the left and right actions of the rings, which I think decategorifies to the anti-involution obtained by composing the bar involution with the KL anti-involution $\omega$.
Also, $\ch(IC_{\omega_1)}$ canonically lives (via pulling back to $G_{KM}$, taking hypercohomology, and then decategoirfying) in $\Hecke{\hh{-1}}{\hh{0}}$ or $\Hecke{\hh{0}}{\hh{1}}$ which via the above anti-involution gets identified with $\Hecke{\hh{1}}{\hh{0}}$}




\subsection{Comparing the \texorpdfstring{$\zeta$}{zeta}-deformed and \texorpdfstring{$q$}{q}-deformed realizations} \label{appendix:comparisonqz}
In this section, we explicitly compare the deformed affine (Frobenius) realization \cref{defn deformed affine frobenius realization} with the following $q-$deformed version.
We assume $n \geq 3$, though there are similar statements in the $n=2$ case (see \cite{EJY1}).
Throughout this section, $\mc{A}_q=\kk[q^{\pm1}]$, $\mc{A}_\zeta=\kk[\zeta^{\pm1}]$ and $\mc{A}_{\zeta,q}=A_{\Z}\otimes_{\Z}\kk$, where $\kk$ is a commutative domain and $A$ is as in \cref{eq coefficient ring with q and zeta}.
Define $V_q$ as follows, over $\mc{A}_q=\kk[q^{\pm1}]$.
\begin{defn}\cite[Definition 2.4]{EJY1}\label{defn q-deformed affine realization}
    Let $V_q$ be the free module over $\mc{A}_q$ with basis $\set{v_i}_{1\leq i \leq n}$ and $\set{v_i^*}$ the corresponding dual basis.
    Let $$\alpha_i=v_i-v_{i+1},\;\; \alpha_{i}^\vee=v_i^*-v_{i+1}^*$$ for $i \neq 0$, and $$\al_0=qv_n-q^{-1}v_1, \;\;\al_0^\vee=q^{-1}v_n^*-qv_1^*.$$
    Then $W$ acts on $V_q$ as follows:
    for the subgroup $W_{\hh{0}}=S_n$ generated by $\set{s_i: i\neq 0}$, we get the usual permutation action, permuting the basis elements $\set{v_k}$.
    The action of $s_0$ is explicitly given by:
    $$s_0(v_n)=q^{-2}v_1,\;\; s_0(v_1)=q^{2} v_n.$$
\end{defn}

Let $V_\zeta$ denote the deformed affine realization over $\mc{A}_\zeta$, as in \cref{defn-deformed affine realization}. 
We then have:
\begin{lemma} \cite[Lemma 2.12]{EJY1}\label{lemma q-deformed to zeta-deformed affine realization}
    There is an isomorphism of $\mc{A}_{\zeta,q}$-modules
    $$V_q\otimes_{\mc{A}_q}\mc{A}_{\zeta,q} \xrightarrow{\sim} V_{\zeta}\otimes_{\mc{A}_\zeta}\mc{A}_{\zeta,q}: v_k \mapsto \zeta^k x_k, \;\text{$1\leq k \leq n$}.$$
    This isomorphism sends the simple roots $\alpha_k, \; 1\leq k \leq n-1$ in $V_q$ to $\zeta^k\alpha_k$ in the target, and sends $\alpha_0$ in $V_q$ to $q^{-1}\alpha_0$.
    The coroots are scaled by the inverse scalars.
\end{lemma}

Let us explicitly describe the Frobenius structure on $V_q$, which generalizes the Frobenius structure used for the $q$-deformed realization for quantum Satake functor in type $A_2$ in \cite{EQuantumI}.
For $s \in S$, we denote by $\ddq_s$ the corresponding Demazure operator in the $q$-deformed affine Frobenius realization, and for an expression $\un{w}$, we denote by $\ddq_{\un{w}}$ the corresponding Demazure operator.

\begin{defn}\label{defn q-deformed affine Frobenius realization}
The \emph{$q$-deformed affine Frobenius realization} is defined as follows.
\begin{itemize}
\item As in \cref{defn deformed affine frobenius realization}, we fix a system of positive roots $(\Phi_q)_I^+$ (or simply $\Phi_I^+$ when there's no ambiguity in the realization being referred to) for each finitary $I \subset S$.
We only define $\Phi_I^+$ for $I$ connected, since we may take a disjoint over the connected components for a more general $I$.

\item For $I \subset S \setminus \set{0}$, we have the standard set of positive roots from the permutation realization of $S_n=W_{\hh{0}}$.
In particular, for connected $I \subset S \setminus\set{0}$, say for $I=\set{i,i+1, \ldots, k}$, we have the positive roots to be $v_j-v_l$ for $i\leq j < l\leq k+1$.

\item Now consider a connected subset $I$ containing $n$. Then $I=\set{i, i+1, \ldots, n, 1,\ldots,k}.$ 
We then define the positive roots to be 
$$\Phi_I^+=\set{v_j-v_l}_{i\leq j<l\leq n}\sqcup \set{v_j-v_i}_{0<j<l\leq k+1}\sqcup \set{qv_j-q^{-1}v_l}_{i \leq j\leq n,\; 0<l\leq k+1}.$$
\item The $\procopdtq_I^J$, $\ddq_I^J$ for each $J \subset I$ are then defined exactly as in \cref{defn deformed affine frobenius realization}.
\end{itemize}
\end{defn}

Note that via the identification in \cref{lemma q-deformed to zeta-deformed affine realization}, the $\procopdtq_I^J$, $\ddq_I^J$ in $V_q$ are scalar multiples of their respective counterparts in $V_\zeta$.
One may show that this gives a Frobenius hypercube as desired, either using the fact above (using the results for $V_\zeta$ and comparing the Frobenius trace maps), or directly using the techniques that we used to show the statement for $V_\zeta$.
We formalize this in the following lemma.
\begin{lemma}
    Let $R=\Sym_{\mc{A}_q}V_q$ For finitary $I\subset J \subset S$, the extension $R^I \hookrightarrow R^J$ in \cref{defn q-deformed affine Frobenius realization} is a Frobenius extension, and $\ddq_I^J$ is a Frobenius trace, with product-coproduct element $\procopdtq_I^J$.
\end{lemma}

Here is the analogue of \cref{lemma demazure operators computation deformed affine}.

\begin{lemma}\label{lemma q-demazure operators computation}
    For $i \in \Om$, let $\ddq_i$ denote the corresponding Demazure operator on $R=\Sym_{\mc{A}_q}V_q$.     
We then have:
\begin{enumerate}
    \item For $i \in \Om$, $\ddq_i(v_j)=\begin{cases}
        1 &\text{ if $j=i<n$}\\
        -1 &\text{ if $j-1=i<n$}\\
        q^{-1} &\text{ if $j=i=n$}\\
        -q &\text{ if $j-1=i=0$}\\
        0  &\text{ otherwise}
    \end{cases}$
    \item For $i,j \in \Om$ such that $s_is_j=s_js_i$, $\ddq_i \ddq_j= \ddq_j \ddq_i$.
    \item For $i \in \Om$, 
    $$\ddq_{i}\ddq_{i+1}\ddq_{i}=
    \begin{cases}
        \ddq_{i+1}\ddq_i\ddq_{i+1} &\text{ if $1\leq i \leq n-2$}\\
        q\ddq_{i+1}\ddq_i\ddq_{i+1} &\text{ if $i=n-1, n$}
    \end{cases}$$
    % \item For $w \in W$, the Demazure operator $\ddq_{\un{w}}$ associated to a reduced expression $\un{w}$ of $w$ is independent of the particular choice of the reduced expression up to a unit in $\mc{A}_q$ of the form $q^l$ for some $l \in \Z$.
    \end{enumerate}
\end{lemma}

% \begin{remark}
% We no longer have symmetry with respect to the rotation operator (see \cref{subsubsec rotational symmetry for deformed affine realization}), which is why the authors chose to work with the $\zeta$-deformed version in the main body of this paper.
% \end{remark}

% \subsubsection{Comparing the Frobenius structures of $V_q$ and $V_\zeta$.}\label{subsubsec comparing q-deformed and zeta-deformed frobenius structures}

We now compare the Frobenius trace maps for the $q$-deformed realization with those of the $\zeta$-deformed realization.
\begin{lemma}\label{lemma comparing Frobenius structures q vs zeta}
As before, let $\dd_s$ denote the Demazure operators  for the $V_\zeta$, and $\ddq_s$ the operators for $V_q$. After base change to $\mc{A}_{\zeta,q}$, we have:
    \begin{enumerate}
        \item $\ddq_i=\zeta^{-k}\dd_i$ for $1 \leq i \leq n-1$.
        \item $\ddq_n=q\dd_n$.
        \item For $I=\set{i,i+1, \ldots,k} \subset S \setminus \set{0},$ we 
        have $\procopdtq_I= \kappa_I \mu_I$ and $\ddq_I=\kappa_I^{-1}\dd_I$, where
        $$\kappa_I= \displaystyle \prod_{i\leq j \leq k}\zeta^{j(k+1-j)}$$
        \item For $I = \set{i, i+1, \ldots, n, 1, \ldots, k}$, we have
        $\procopdtq_I= \kappa_I \mu_I$ and $\ddq_I=\kappa_I^{-1}\dd_I$, where
        \begin{align}\label{eqn 1 lemma comparing Frobenius structures q vs zeta}
        \notag \kappa_I &= q^{(n+1-i)(k+1)}\left(\prod_{i\leq j \leq n}\zeta^{j(n+k+1-j)}\right)\left( \prod_{1\leq j\leq k}\zeta^{j(k+1-j)}\right)\\
        &=q^{(n+1+k-i)(k+1)}\left(\prod_{i \leq j \leq n+k} \zeta^{j(n+k+1-j)}\right).
        \end{align}
        \item For an arbitrary finitary subset $I \subset S$, let $I_1, \ldots , I_k$ be the connected components of $I$.
        Then $\procopdtq_I=\kappa_I\mu_I$ and $\ddq_I=\kappa_I^{-1}\mu_I$, where $$\kappa_I=\kappa_{I_1}\kappa_{I_2}\ldots \kappa_{I_k}.$$
        \item For $J\subset I \subset S$ finitary subsets, we have $\procopdtq_I^J= \kappa_I^J \mu_I^J$ and $\ddq_I^J=(\kappa_I^J)^{-1}\dd_I^J$ where $\kappa_I^J=\kappa_J^{-1}\kappa_I.$
    \end{enumerate}
\end{lemma}
\begin{proof}
These are all straightforward computations.
\end{proof}

%     (1) and (2) are straightforward computations.
%     For (3), observe that
%     \begin{align*}
%         \procopdtq_I = \prod_{i\leq j <l\leq k+1}(v_j-v_k)= \prod_{i\leq j<l\leq k+1}(\zeta^jx_j-\zeta^lx_l) &= \prod_{i\leq j\leq k}\zeta^{j(k+1-j)}\prod_{j<l\leq k+1}(x_j-\zeta^{l-j}x_l)\\
%         &=\left(\prod_{i\leq j\leq k}\zeta^{j(k+1-j)}\right)\mu_I.
%     \end{align*}
%     The statement regarding $\ddq_I$ follows.\\
%     For (4), we have
%     \begin{align*}
%         \procopdtq_I &= \left(\prod_{i\leq j<l \leq n}(v_j-v_l)\right)\left( \prod_{1\leq j<l<k+1}(v_j-v_l)\right)\left(\prod_{\tiny\begin{array}{c} i\leq j \leq n\\ 1\leq l \leq k+1 \end{array}}q(v_j-q^{-2}v_l) \right)\\
%         &=\left(\prod_{i\leq j<l \leq n}\zeta^j(x_j-\zeta^{l-j}x_l)\right)\left( \prod_{1\leq j<l<k+1}\zeta^j(x_j-\zeta^{l-j}x_l)\right)\left(\prod_{\tiny\begin{array}{c} i\leq j \leq n\\ 1\leq l \leq k+1 \end{array}}q\zeta^j(x_j-\zeta^{n+l-j}x_l) \right)\\
%         &=\left(\prod_{i\leq j \leq n}\zeta^{j(n-j)}\right)\left( \prod_{1\leq j\leq k}\zeta^{j(k+1-j)}\right)\left(\prod_{ i\leq j \leq n} q^{k+1}\zeta^{j(k+1)}\right) \mu_I\\
%         &=q^{(n+1-i)(k+1)}\left(\prod_{i\leq j \leq n}\zeta^{j(n+k+1-j)}\right)\left( \prod_{1\leq j\leq k}\zeta^{j(k+1-j)}\right) \mu_I.
%     \end{align*}
%     The last equality in \cref{eqn 1 lemma comparing Frobenius structures q vs zeta} follows from the equality:
%     \begin{align*}\prod_{1 \leq j \leq k} \zeta^{j(k+1-j)}=\prod_{n+1 \leq j \leq n+k}\zeta^{(j-n)(n+k+1-j)} &= \prod_{n+1 \leq j \leq n+k}\zeta^{j(n+k+1-j)}q^{2(n+k+1-j)}\\
%     &=q^{k(k+1)}\prod_{n+1 \leq j \leq n+k}\zeta^{j(n+k+1-j)}.
%     \end{align*}
%     The statement for $\ddq_I$ follows.\\
%     (5) follows from how $\procopdtq_I, \ddq_I$ are defined for disconnected finitary subsets $I$, and (6) follows from the fact that $\procopdtq_I^J\procopdtq_J=\procopdtq_I$, $\ddq_I^J\ddq_J=\ddq_I$ (and the analogous formulas for the $\mu_I^J$ and $\dd_I^J$).
% \end{proof}


% Associated to the Frobenius hypercube from the $q$-deformed affine realization, we have a diagrammatic category generated by oriented cups, caps and oriented crossings, as in \cref{subsec-diagrammatics for frob hypercube}, and we may construct $\qdiag$ which is a subcategory of $\SSBim_q$ constructed from $V_q$. 
% We use $\zdiag$ to refer to $\DiagSBS$ from \cref{subsec-diagrammatic singular bott-samelson bimodules} . 
% We now explicitly describe how the generating diagrams for $V_q$ compare with those of $V_\zeta$ (inside the algebraic category of singular Soergel bimodules). We work over $\mc{A}_{\zeta,q}$, so that we have $\SSBim_q$
% This will be essential to produce a $q-$form $\GS_q$ for the $\GS_\zeta$ that we have constructed. \BE{I strongly prefer to phrase this as an equivalence from $\zdiag$ to $\qdiag$ which gives a commuting square involving: the functors to $SSBim_q$ and $SSBim_{\zeta}$ and an equivalence between the SSBims coming from that lemma above.}

Let $\qdiag$ denote the diagrammatic category associated to the Frobenius hypercube from \cref{defn q-deformed affine Frobenius realization}, and $\zdiag$ for the one associated to \cref{defn deformed affine frobenius realization} (see \cref{subsec-diagrammatic singular bott-samelson bimodules}).
Let $\Phi_q: \qdiag \rightarrow \SBSBim_q$, $\Phi_\zeta: \zdiag \rightarrow \SBSBim_\zeta$ be the corresponding functors to the algebraic categories associated to $V_q$ and $V_\zeta$ respectively. Over $\mc{A}_{q,\zeta},$ we have an equivalence of 2-categories $\Theta^\zeta_q:\SSBim_q\xrightarrow{\sim}\SSBim_\zeta$ using \cref{lemma q-deformed to zeta-deformed affine realization}.
Working over $\mc{A}_{q,\zeta}$, we then have

\begin{equation}
\begin{tikzcd}
    \qdiag \arrow[r, hook, "\Phi_q"] \arrow[d, "\sim" sloped, "\vartheta_q^\zeta"'] & \SSBim_q \arrow[d, "\sim" sloped , "\Theta^\zeta_q"'] \\
    \zdiag \arrow[r, hook, "\Phi_\zeta"] & \SSBim_\zeta
\end{tikzcd}
\end{equation}
where $\vartheta$ is explicitly described below on the generators of $\qdiag$. This will be essential to produce a $q$-form $\GS_q$ for the $\GS_\zeta$ that we have constructed.

Recall from \cref{subsec-diagrammatics for frob hypercube} that the clockwise cap and counter-clockwise cup are mapped to the multiplication maps and inclusion maps for the respective rings, and is independent of the particular choice of the Frobenius structure we use.
Hence $\vartheta$ sends the clockwise caps (resp. counter-clockwise cups) in $\qdiag$ to those in $\zdiag$.

The counter-clockwise caps and clockwise cups are sent to the corresponding Frobenius trace maps and the co-product maps respectively in the algebraic category. 
The images of these diagrams in $\SSBim_q \simeq \SSBim_\zeta$ then differ by scalars specified by the $\kappa_I$ as in \cref{lemma comparing Frobenius structures q vs zeta}.
More precisely, we have:

\begin{equation}\label{diag- comparing q cap cc with zeta cap cc}
\vartheta^\zeta_q 
\left( \;
%begin anticlockwise cup
{\begin{tikzpicture}[baseline=(current bounding box.center),
    scale=1,
    line width=0.8pt,
    % Updated style to place two arrows on each line
    arrow_style/.style={
        decoration={
            markings,
            mark=at position 0.3 with {\arrow{stealth[black]}}, % First arrow
            mark=at position 0.8 with {\arrow{stealth[black]}}  % Second arrow
        },
        postaction={decorate}
    }
]
\draw[color={rgb,255:red,0; green,0; blue,192}, arrow_style] (0.8,0) arc (0:180:0.8);
\node at (0,1.5) {$\gr{I}$};
\node at (0,0.3) {$\gr{I\setminus \set{s}}$};
\node at (1,1) {}; %dummy node
\node at (-1,1) {}; %dummy node
\node at (0,1.9) {}; %dummy node

%bounding box
\node[draw, inner sep=0pt, fit=(current bounding box)] {};
\end{tikzpicture}_{\qdiag}}
\right)
%end  anticlockwise cap
= (\kappa_I^{I \setminus \set{s}})^{-1}
\;
%begin anticlockwise cap
{\begin{tikzpicture}[baseline=(current bounding box.center),
    scale=1,
    line width=0.8pt,
    % Updated style to place two arrows on each line
    arrow_style/.style={
        decoration={
            markings,
            mark=at position 0.3 with {\arrow{stealth[black]}}, % First arrow
            mark=at position 0.8 with {\arrow{stealth[black]}}  % Second arrow
        },
        postaction={decorate}
    }
]
\draw[color={rgb,255:red,0; green,0; blue,192}, arrow_style] (0.8,0) arc (0:180:0.8);
\node at (0,1.5) {$\gr{I}$};
\node at (0,0.3) {$\gr{I\setminus \set{s}}$};
\node at (1,1) {}; %dummy node
\node at (-1,1) {}; %dummy node
\node at (0,1.9) {}; %dummy node

%bounding box
\node[draw, inner sep=0pt, fit=(current bounding box)] {};
\end{tikzpicture}_{\zdiag}}
%end  anticlockwise cap
\end{equation}

\begin{equation}\label{diag- comparing q cup c with zeta cup c}
\vartheta^\zeta_q 
\left( \;
%begin anticlockwise cup
{\begin{tikzpicture}[baseline=(current bounding box.center),
    scale=1,
    line width=0.8pt,
    % Updated style to place two arrows on each line
    arrow_style/.style={
        decoration={
            markings,
            mark=at position 0.3 with {\arrow{stealth[black]}}, % First arrow
            mark=at position 0.8 with {\arrow{stealth[black]}}  % Second arrow
        },
        postaction={decorate}
    }
]
\draw[color={rgb,255:red,0; green,0; blue,192}, arrow_style] (0.8,1.9) arc (360:180:0.8);
\node at (0,1.5) {$\gr{I}$};
\node at (0,0.4) {$\gr{I\setminus \set{s}}$};
\node at (1,0) {}; %dummy node
\node at (-1,0) {}; %dummy node
%bounding box
\node[draw, inner sep=0pt, fit=(current bounding box)] {};
\end{tikzpicture}_{\qdiag}}
%end  anticlockwise cap
\right)
= \kappa_I^{I \setminus \set{s}}
\;
%begin anticlockwise cup
{\begin{tikzpicture}[baseline=(current bounding box.center),
    scale=1,
    line width=0.8pt,
    % Updated style to place two arrows on each line
    arrow_style/.style={
        decoration={
            markings,
            mark=at position 0.3 with {\arrow{stealth[black]}}, % First arrow
            mark=at position 0.8 with {\arrow{stealth[black]}}  % Second arrow
        },
        postaction={decorate}
    }
]
\draw[color={rgb,255:red,0; green,0; blue,192}, arrow_style] (0.8,1.9) arc (360:180:0.8);
\node at (0,1.5) {$\gr{I}$};
\node at (0,0.4) {$\gr{I\setminus \set{s}}$};
\node at (1,0) {}; %dummy node
\node at (-1,0) {}; %dummy node
%bounding box
\node[draw, inner sep=0pt, fit=(current bounding box)] {};
\end{tikzpicture}_{\zdiag}} \quad .
%end  anticlockwise cap
\end{equation}
\vspace{5 pt}


The generating upward and downward crossings (see \cref{diag generating crossings}) in both $\qdiag$ and $\zdiag$ go to the exact same 2-morphism in $\SBSBim_q \simeq \SBSBim_\zeta$ so that $\vartheta^\zeta_q$ sends these crossings in $\qdiag$ to those in $\zdiag$ (without any coefficients).

\begin{remark}
From here one could explicitly describe a functor $\GS_q': \cwebs \rightarrow \qdiag$ over $\mc{A}_{q, \zeta}$ when $\kk=\Q$, by composing $\GS_{\zeta}$ with $(\vartheta_q^\zeta)^{-1}$. However, this functor does not deserve to be called $\GS_q$, as it is not defined over $\mc{A}_q$ but only over an extension. 
One needs to make a careful, non-symmetric choice of the scalars $\lambda$ and $\nu$ from \cref{diag-dQGS on 2-morphisms} in order for the result to land in $\qdiag$ (over $\mc{A}_q$). We hope to provide this choice in an updated version of the paper. \PTv2{Add it.}
\end{remark}



% Using how the generating 2-morphisms compare between $\qdiag$ and $\zdiag$ after passing to $\SBSBim$, we can now compare the images of any diagrams from $\qdiag$ and $\zdiag$.

% Specifically for the purposes of constructing $\GS_q$, we now explicitly compare the diagrams in the RHS of \cref{diag-dQGS on 2-morphisms} (which are in $\zdiag$) with the corresponding diagrams in $\qdiag$.
% We shall now drop writing the $\Phi_q, \Phi_\zeta$ when comparing diagrams for notational simplicity.
% Any equality between diagrams in $\qdiag$ and $\zdiag$ should be thought of as an equality only after passing to the algebraic category.

% \begin{align}
%      \isvg{0.15}{cap_ssbim_c}_{\qdiag} &=  \isvg{0.15}{cap_ssbim_c}_{\zdiag}(\kappa_{\hh{a}}^{\hh{a,a-k}})^{-1}\\
%      \hspace{20 pt}\isvg{0.15}{cap_ssbim_cc}_{\qdiag} &= \isvg{0.15}{cap_ssbim_cc}_{\zdiag} (\kappa_{\hh{a}}^{\hh{a,a+k}})^{-1}
%      \\
%      \isvg{0.15}{cup_ssbim_c}_{\qdiag} &= \isvg{0.15}{cup_ssbim_c}_{\zdiag}\kappa_{\hh{a}}^{\hh{a,a-k}}\\
%      \hspace{20 pt}\isvg{0.15}{cup_ssbim_cc}_{\qdiag} &= \isvg{0.15}{cup_ssbim_cc}_{\zdiag}\kappa_{\hh{a}}^{\hh{a,a+k}}\\
%      \isvg{0.15}{image_of_up_trivalent}_{\qdiag}  &= \isvg{0.15}{image_of_up_trivalent}_{\zdiag}\kappa_{\hh{\small a,a+k+l}}^{\hh{\small a,a+l, a+k+l}}\\
%      \isvg{0.15}{image of down trivalent}_{\qdiag} &= \isvg{0.15}{image_of_down_trivalent}_{\zdiag}
% \end{align}

% We now explicitly compute all the coefficients appearing in the above equalities.
% \begin{lemma} \label{lemma q to zeta frobenius coefficient computations}
% Let $0 \leq a <n, 1 \leq k \leq n$.
%     \begin{enumerate}
%         \item $\displaystyle\kappa^{\hh{a,k}}_{\hh{a}}= q^{a(k-a)}\prod_{a+1\leq j \leq k}\zeta^{j(n+a-k)}$ if $a \leq k \leq n.$
%         \item $\displaystyle\kappa^{\hh{a,k}}_{\hh{a}}= q^{(n-a)(a-k)}\left(\prod_{a+1\leq j \leq n}\zeta^{j(a-k)}\right)\left( \prod_{1\leq j \leq k}\zeta^{j(a-k)}\right)$ if $k \leq a < n$.
%         \item $\displaystyle\kappa^{\hh{a,a+k}}_{\hh{a}}= q^{ak}\prod_{a+1\leq j \leq k+a}\zeta^{j(n-k)}$ if $a+k \leq n.$
%         \item $\displaystyle\kappa^{\hh{a,a+k}}_{\hh{a}}= q^{(n-a)(n-k)}\left(\prod_{a+1\leq j \leq n}\zeta^{j(n-k)}\right)\left( \prod_{1\leq j \leq a+k-n}\zeta^{j(n-k)}\right)$ if $a+k \geq n.$
%         \item $\displaystyle\kappa^{\hh{a,a-k}}_{\hh{a}}= q^{a(n-k)}\prod_{a+1\leq j \leq a-k+n}\zeta^{jk}$ if $a-k \leq 0.$
%         \item $\displaystyle\kappa^{\hh{a, a-k}}_{\hh{a}}= q^{(n-a)k}\left(\prod_{a+1\leq j \leq n}\zeta^{jk}\right)\left( \prod_{1\leq j \leq a-k}\zeta^{jk}\right)$ if $a-k \geq 0$.
%         \item $\displaystyle\kappa_{\hh{a,a+k+l}}^{\hh{a,a+l,a+k+l}}=\prod_{a+1\leq j \leq a+l}\zeta^{jk}$ for $a+k+l-1 < n$.
%         \item $\displaystyle\kappa_{\hh{a,a+k+l}}^{\hh{a,a+l,a+k+l}}=q^{l(a+k+l-n)} \prod_{a+1\leq j \leq a+l}\zeta^{jk}$ for $a+k+l-1 \geq n$, $a+l \leq n$.
%         \item $\displaystyle\kappa_{\hh{a,a+k+l}}^{\hh{a,a+l,a+k+l}}=q^{(n-a)k} \left(\prod_{a+1\leq j \leq n}\zeta^{jk}\right)\left(\prod_{1\leq j \leq a+l-n}\zeta^{jk}\right)$ for $a+k+l-1 > n$, $a+l > n$.
%     \end{enumerate}
% \end{lemma}
% \begin{proof}
%     These are direct computations using \cref{lemma comparing Frobenius structures q vs zeta} which are left as an exercise to the reader.
% \end{proof}

%\subsubsection{From $\GS_\zeta$ to $\GS_q$.}

%\input{For future/Add to supplementary}
%add a detailed discussion going from usual geometric Satake to mSSBim for affine Kac Moody realization.
%explain how to think of $\SSBim_q$ or $\SSBim_\zeta$ as a formal deformation of $SSBim$ for the Kac-Moody realization.


\subsection{Symmetries of Diagrammatic and Algebraic Bott-Samelson bimodules.}\label{subsec reversal duality and phi}

This subsection is dedicated to explaining precisely how the diagrammatic symmetries of horizontal and vertical flipping (with suitable arrow reversals) connect to certain standard symmetries of $\SBSBim$ (and hence $\SSBim$). Nothing in this section should surprise the experts.
While we only discuss this in the context of $\SBSBim$ from the deformed affine Frealization, the contents of this section generalizes to any (partial, graded or ungraded) Frobenius hypercube $\Gamma$, with $\SBSBim$ replaced by the algebraic category $C(\Gamma)$ from \cite{EWSFrob}.

We retain notation from \cref{subsec-diagrammatic singular bott-samelson bimodules}.


\begin{defn}\label{defn reverse of singular multistep expression}
  Given a singular multistep expression  
  $$IK_{\bullet}=[[I_0 \subset K_1 \supset I_1 \subset \ldots \subset K_{d} \supset I_d]],$$
  we define the \emph{reverse of the multistep expression} $IK_\bullet$ to be
  $$\Rev(IK_\bullet):=[[I_d \subset K_{d} \supset I_{d-1} \subset \ldots \supset I_1 \subset K_1 \supset I_0 ]].$$
\end{defn}

\begin{defn}
    Let $\Rev$ denote the 2-functor on $\mc{D}^{\pre}$ which fixes regions, flips diagrams horizontally (i.e., reflects diagrams across a vertical axis) and  fixes floating polynomials. This 2-functor also necessarily changes some strand orientations to make the diagram consistent with the region labels. 
    Note that for a 1-morphism in $\mc{D}^{\pre}$ given by a singlestep expression $I_\bullet$, $\Rev(I_\bullet)$ is just the reverse of the singlestep expression $I_\bullet$.
\end{defn}

\begin{defn}
    Let $\Dual$ denote the 2-functor on $\mc{D}^{\pre}$ which fixes regions and flips diagrams vertically (i.e., reflects diagrams across a horizontal axis) and fixes floating polynomials. This 2-functor also necessarily changes some strand orientations to make the diagram consistent with the region labels. Note that for a 1-morphism in $\mc{D}^{\pre}$ given by a singlestep expression $I_\bullet$, $\Dual(I_\bullet)=I_\bullet$.
\end{defn}

There are a few conventions in the literature for the corresponding functors on $\SBSBim$, we follow those in \cite[(3.26), (3.48)]{BBE}.     
 

\begin{defn}\label{def reversal symmetry R algebraic}
    Let $\Rev$ be the 2-functor on $\SBSBim$ defined as follows.
    $\Rev$ fixes objects, and acts on the 1-morphism categories by
    $$\Hom_{\SBSBim}(I,J) \xrightarrow{\sim} \Hom_{\SBSBim}(J,I): \Bimod{I}{J}{M} \mapsto \Bimod{J}{I}{M}(\ell(J)-\ell(I)).$$
    That is, $\Rev$ swaps the module actions on the left and right (using commutativity of the rings $R^I, R^J$) and adds a grading shift; $\Rev$ does not change the underlying functions of $2$-morphisms.
    This functor is a weak 2-functor which is contravariant on 1-morphisms and covariant on 2-morphisms, and as such includes the coherence data given by the 2-morphisms
    \begin{gather}
       \Rev(\Bimod{R^J}{R^K}{N})\otimes_{R^J}\Rev(\Bimod{R^I}{R^J}{M}) 
       \xrightarrow{\sim}\Rev(\Bimod{R^I}{R^J}{M} \otimes_{R^J} \Bimod{R^J}{R^K}{N}): n\otimes m \mapsto m\otimes n\\
       \Rev(\Bimod{R^I}{R^I}{R^I})=\Bimod{R^I}{R^I}{R^I} \xrightarrow[\id]{\sim} \Bimod{R^I}{R^I}{R^I}
    \end{gather}
    for $\Bimod{R^I}{R^J}{M} \in \Hom_{\SBSBim}(J,I)$, $\Bimod{R^J}{R^K}{N} \in \Hom_{\SBSBim}(K,J)$, $I,J\subset S$ finitary.
\end{defn}

\begin{defn}\label{def duality symmetry D algebraic}
    Let $\Dual$ denote the 2-functor on $\SBSBim$ defined as follows. 
    $\Dual$ fixes objects, and acts on the 1-morphism categories by
    $$\Hom(I,J)^{\op}\to \Hom(I,J): M \mapsto \Hom^\bullet_{(-,R^J)}(M, R^J)(2\ell(I)-2\ell(J)),$$ 
     where $\Hom^\bullet_{(-,R^J)}(B, R^J)$ denotes the space of right graded $R^J$-module homomorphisms into $R^J$, with $((r \otimes s)\cdot f)(v)=s(f(rv))$ for $r \in R^I$, $s\in R^J$, $f \in B^\vee$.
    The 2-functor $\Dual$ is a weak 2-functor which is covariant on 1-morphisms and contravariant on 2-morphisms, and as such includes the coherence data given by the 2-morphisms 
    \begin{gather}
    \Dual(M)\otimes_{R^J} \Dual(N) \xrightarrow[\epsilon_{M,N}]{\sim} \Dual(M\otimes_{R^J} N): f\otimes g \mapsto \bigl( m \otimes n \mapsto g(f(m)\cdot n) \bigr),\label{eq coherence data for duality}\\
    \Dual(\Bimod{J}{J}{R^J}) \xrightarrow[\varepsilon_J]{\sim} \Bimod{J}{J}{R^J}: f \mapsto f(1), \label{eq coherence data duality identity 1-morphism}
    \end{gather}
    for $M \in \Hom_{\SBSBim}(J,I), N \in \Hom_{\SBSBim}(K,J)$ for finitary $I,J \subset S$. 
    Note that the map in \cref{eq coherence data for duality} is an isomorphism because $M$, being a 1-morphism in $\SBSBim$, is a finitely generated projective (in fact, free) graded right $R^J$-module.
    We refer to $\Dual$ as the \emph{(right) duality functor} on $\SBSBim$.
\end{defn}


\begin{remark}
    By definition, the 2-functors $\Rev$ and $\Dual$ defined on $\mc{D}^{\pre}$ and $\SBSBim$ preserve objects. As such, natural transformations between 2-functors appearing in this section will only involve the data of 2-morphisms satisfying certain compatibility conditions (we default to the identity 1-morphism at the object level).
\end{remark}

\begin{remark}\label{rmk sbsbim larger for weel definedness of dual and rev}
    Strictly speaking, for $\Rev$ and $\Dual$ to be well-defined, we should allow $\SBSBim$ to be slightly bigger then (but equivalent to) the one in \cref{def sbsbim}: namely, we should allow 1-morphism categories to be closed under isomorphisms in the corresponding category of graded bimodules, so that 1-morphisms are those isomorphic (not just equal) to singular Bott-Samelson bimodules. That this category is closed under $\Rev$ and $\Dual$ then follows from \cref{lemma self duality of Bott-Samelsons} and \cref{lemma reversal on bott samelsons and multistep expressions} below, or by direct checks on the monoidal generators.
\end{remark}


The following lemma is easy and left as an exercise to the reader.
\begin{lemma}\label{lemma dual and rev are involutions}
    The 2-functors $\Dual$ and $\Rev$ on $\mc{D}^{\pre}$ and $\SBSBim$ are involutions.
    More precisely, there are equalities
    $$\Dual \circ \Dual = \Id_{\mc{D}^{\pre}}= \Rev \circ \Rev$$
    of 2-functors on $\mc{D}$, and an equality and an isomorphism
    $$\Id_{\SBSBim}=\Rev\circ \Rev, \qquad \Id_{\SBSBim} \Rightarrow \Dual \circ \Dual$$ of 2-functors on $\SBSBim$, where the latter is given by the evaluation maps $$B \xrightarrow{\sim} \Dual\circ \Dual(B): m \mapsto(f\mapsto f(m))$$ for 1-morphisms $B$ in $\SBSBim$.
\end{lemma}

Recall that 1-morphisms in $\SBSBim$ are\footnote{up to an isomorphism, see \cref{rmk sbsbim larger for weel definedness of dual and rev}.} given by singular Bott-Samelson bimodules
\begin{equation*}
    \BS(IK_{\bullet}):= R^{I_0}\otimes_{R^{K_1}} R^{I_1}\otimes_{R^{K_2}} \cdots \otimes_{R^{K_d}}R^{I_d}(\sum_{i=1}^d \ell(K_i)-\ell(I_i)) \in \Hom_{\SBSBim}(I_d,I_0),
\end{equation*}
for a singular multistep expression
\begin{equation}\label{eq singular multistep expression supplementary}
    IK_{\bullet}=[[I_0 \subset K_1 \supset I_1 \subset \ldots \subset K_{d} \supset I_d]].
\end{equation}

The following lemma is immediate from the definitions.
\begin{lemma}\label{lemma reversal on bott samelsons and multistep expressions}
    For any  multistep expression $IK_\bullet$ as above, there is an isomorphism 
    $$\etadot_{IK_\bullet}: \BS(\Rev(IK_\bullet)) \xrightarrow{\sim} \Rev(\BS(IK_\bullet)):r_d\otimes \cdots \otimes r_0 \mapsto r_0 \otimes \cdots \otimes r_d$$ in $\Hom_{\SBSBim}(I_0, I_d)$, for $r_j \in R^{I_j}$, $0 \leq j \leq d$.
\end{lemma}

The following lemma uses Frobenius structures to relate Bott-Samelson bimodules with their duals.

\begin{lemma}\label{lemma self duality of Bott-Samelsons}
    Fix a singular multistep expression $IK_\bullet$ as in \cref{eq singular multistep expression supplementary}. 
    The map 
    $$\eta_{IK_\bullet}: \BS(IK_{\bullet}) \to \Hom^\bullet_{R^{I_d}}(\BS(IK_{\bullet}), R^{I_d})(2(\ell(I_0)-\ell(I_d)))=D(\BS(IK_\bullet))$$ 
    given by
    \begin{equation}
        r_0 \otimes \dots \otimes r_d \mapsto \left(s_0\otimes \cdots \otimes s_d \mapsto \dd^{I_{d-1}}_{K_d}\left(\cdots \dd^{I_1}_{K_2}(\dd^{I_0}_{K_1}(r_0s_0)r_1s_1) \cdots r_{d-1}s_{d-1} \right)r_ds_d \right) \label{eq self duality of bott samelsons}
    \end{equation}
   for $r_j, s_j \in R^{I_j}$, $0 \leq j \leq d$, is an isomorphism in $\Hom_{\SBSBim}(I_d,I_0)$.
   Hence singular Bott-Samelson bimodules are self-dual (canonically, after choosing Frobenius structures). 
\end{lemma}
\begin{proof}
    The given map is a homomorphism of graded $(R^{I_d}, R^{I_0})$-bimodules, so we just have to show that the given map is an isomorphism of sets. We suppress the grading shifts and grading for notational simplicity.
    
    We proceed by induction on $d$. For $d=0$, the result is trivial.
    Now suppose the map 
    \begin{gather}
      \notag \eta': R^{I_0}\otimes_{R^{K_1}} \cdots \otimes_{R^{K_{d-1}}}R^{I_{d-1}} \rightarrow  \Hom^\bullet_{R^{I_{d-1}}}(R^{I_0}\otimes_{R^{K_1}} \cdots \otimes_{R^{K_{d-1}}}R^{I_{d-1}}, R^{I_{d-1}})\\
      \notag  r_0 \otimes \cdots \otimes r_d 
      \mapsto (s_0\otimes \cdots \otimes s_d \mapsto \dd^{I_{d-2}}_{K_{d-1}}\left(\cdots \dd^{I_1}_{K_2}(\dd^{I_0}_{K_1}(r_0s_0)r_1s_1) \cdots r_{d-2}s_{d-2} \right)r_{d-1}s_{d-1})
    \end{gather}
    is an isomorphism.
    Let us denote $ R^{I_0}\otimes_{R^{K_1}} R^{I_1}\otimes_{R^{K_2}} \cdots \otimes_{R^{K_{d-1}}}R^{I_{d-1}}$ by $M$.
     The map that we'd like to show is an isomorphism is then
    \begin{gather*}
       \eta_{IK_\bullet}:M\otimes_{R^{K_d}}R^{I_d} \rightarrow \Hom^\bullet_{R^{I_d}}(M \otimes_R^{K_d}R^{I_d}, R^{I_d}) \\
       m \otimes r \mapsto \left(m' \otimes s \mapsto \dd^{I_{d-1}}_{K_d}\left((\eta'(m))(m')\right)r s\right).
    \end{gather*}
The map $\eta_{IK_\bullet}$ may be factored as

\[
\begin{tikzcd}
   M \otimes_{R^{K_d}} R^{I_d} \arrow[r, "\eta_{IK_\bullet}"] \arrow[d, "\eta'\otimes \id"] &\Hom^\bullet_{R^{I_d}}(M \otimes_{R^{K_d}} R^{I_d}, R^{I_d}) \\
   \Hom^\bullet_{R^{I_{d-1}}}(M, R^{I_{d-1}})\otimes_{R^{K_d}} R^{I_d} \arrow[r, "\upsilon \otimes \id"] & \Hom^\bullet_{R^{K_d}}(M ,R^{K_d})\otimes_{R^{K_d}} R^{I_d} \arrow[u, "\beta"],
\end{tikzcd}\]
    where
    \begin{align*}
        \upsilon: f\mapsto  (m \mapsto \dd^{I_{d-1}}_{K_d}(f(m))) \qquad &f \in \Hom^\bullet_{R^{I_{d-1}}}(M,R^{I_{d-1}}),\; m \in M\\
        \beta: g \otimes r \mapsto (m'\otimes s\mapsto g(m')rs) \qquad&g \in \Hom^\bullet_{R^{K_d}}(M,R^{K_d}), \; m' \in M,\;\; r,s \in R^{I_d}.
    \end{align*}
    Both $\upsilon$ and $\beta$ are isomorphisms since $M$ is a finitely generated projective (in fact, free) graded module over the rings $R^{I_d}$ and $R^{K_d}$; $\eta'$ is an isomorphism by induction hypothesis.
    It follows that $\eta_{IK_\bullet}$ is an isomorphism as desired.
\end{proof}


Recall that we had the 2-functor $\Phi: \mc{D}^{\pre} \rightarrow \SBSBim$ which sends diagrams to corresponding morphisms between bimodules.
\begin{theorem}\label{thm phi commutes with dual and rev}
    The 2-functor $\Phi$ intertwines the 2-functors $\Rev$ and $\Dual$ on $\mc{D}^{\pre}$ with those on $\SBSBim$. 
    More precisely, there are isomorphisms of 2-functors 
    \begin{equation}
        \Phi \circ \Rev \xrightarrow{\sim} \Rev \circ \Phi, 
        \qquad
        \Phi \circ \Dual \xrightarrow{\sim} \Dual \circ \Phi \label{eq phi commutes with rev and dual}
    \end{equation} 
    given by the data
    \begin{gather}
    \etadot_{I_\bullet}:\Phi \circ \Rev (I_\bullet)=\BS(\Rev(I_\bullet)) \xrightarrow{\sim} \Rev(\BS(I_\bullet))= \Rev\circ \Phi(I_\bullet)
    \\
    \eta_{I_\bullet}:\Phi \circ \Dual (I_\bullet)=\BS(I_\bullet) \xrightarrow{\sim} \Dual(\BS(I_\bullet))=\Dual \circ \Phi (I_\bullet)
    \end{gather}
    for singlestep expressions $I_\bullet$, where $\etadot$ and $\eta$  are as in  \cref{lemma reversal on bott samelsons and multistep expressions} and \cref{lemma self duality of Bott-Samelsons}.
\end{theorem}
\begin{proof}
    Checking \cref{eq phi commutes with rev and dual} involves checking the compatibility of $\eta$ and $\etadot$ with the coherence 2-morphisms in \cref{def duality symmetry D algebraic} and \cref{def reversal symmetry R algebraic} respectively at the 1-morphism level, as well as checking certain equalities on (generating) 2-morphisms.
    Both are routine checks, and follows from the explicit formulas for the coherence data in \cref{def reversal symmetry R algebraic}, \cref{def duality symmetry D algebraic}, and the formulas for the isomorphisms in \cref{lemma self duality of Bott-Samelsons} and \cref{lemma reversal on bott samelsons and multistep expressions}.
    
    Here's an example to illustrate the check on the generating 2-morphisms. To check \cref{eq phi commutes with rev and dual} for $\Dual$ on a clockwise cap with region labels $I$ and $Ii$, we need to show that the diagram
    \begin{equation*}
    \begin{tikzcd}
        R^I \arrow[r, "\eta_{[I]}"] \arrow[d, "\Delta"] & \Dual(R^I) \arrow[d, "\Dual(m)"] \\
        R^I \otimes_{R^{Ii}} R^I(\ell(Ii)-\ell(I)) \arrow[r, "\eta_{[I\subset Ii \supset I]}"] &\Dual\left(R^I\otimes_{R^{Ii}}R^I(\ell(Ii)-\ell(I))\right)
    \end{tikzcd}            
    \end{equation*}
    commutes. 
    This is straightforward to check; we have
    \begin{equation*}
        \Dual(m)\circ \eta_{[I]}: r \mapsto (s\mapsto sr) \mapsto \bigl( s_1 \otimes s_2 \mapsto s_1 s_2 r \bigr)
    \end{equation*}
    which agrees with
    \begin{equation}
      \eta_{[I \subset Ii \supset I]}\circ \Delta: r\mapsto \sum_i e_i \otimes f_i r\mapsto \left(s_1 \otimes s_2 \mapsto \sum_i\dd^I_{Ii}(e_is_1)f_i rs_2 = s_1 s_2 r\right)  
    \end{equation}
    where $\{e_i\}, \{f_i\}$ are dual bases for $R^{Ii}\hookrightarrow R^I$ (so that $\sum_i\dd^{I}_{Ii}(se_i)f_i=s$ for any $s \in R^I$).
\end{proof}

\begin{defn}\label{def rotation by 180 degrees}
    Let $\rotation$ be the 2-functor on $\mc{D}^{\pre}$, given by rotating a diagram by 180 degrees. This is a strict 2-functor which coincides with the composition $\Rev \circ \Dual=\Dual \circ \Rev$.
    On $\SBSBim$, we define $\rotation$ to be the (weak) 2-functor given by the composition $\Dual \circ \Rev$.
\end{defn}


The following then easily follows from \cref{thm phi commutes with dual and rev}.
\begin{corollary}\label{corollary phi intertwines rotation}
    The 2-functor $\Phi: \mc{D}^{\pre} \rightarrow \SBSBim$ intertwines $\rotation$ with the 2-functor $\rotation= \Dual \circ \Rev$ (as well as with the 2-functor $\Rev \circ \Dual$) on $\SBSBim$. 
\end{corollary}

%    \BE{please add some sentence like this: In particular, $\rotation$ agrees with the isomorphism coming from adjunction, where the units and counits of adjunction come from the Frobenius structure....} \BE{Proof: because it works diagrammatically... then maybe remark somewhere, not necessarily in the proof, maybe before the statement, that this isn't surprising given that the frobenius structure is used to build the map $\eta$.}\PT{I'm not sure about ``because it works diagrammatically", unless we have 2-fullness for $\Phi$. The statement you want to make still holds, purely by induction restriction biadjunctions, but I think it is better not added as part of the corollary. Adding it as a remark below.} \BE{I think you're emphasizing the wrong point here.}\PT{Was a work in progress. Will wait for you to finish.} \BE{finished, running to the next meeting...}\PT{Great, let's maybe meet some time over Zoom tomorrow and post the paper to arxiv (after one last read from beginning to end)?}

\begin{remark}\label{rmk rotation and flips and adjunctions}
In $\mc{D}^{\pre}$, the rotation $2$-functor $\rotation$ is more than just the composition of two flip symmetries; it is also realized by composition with caps and cups. That is, $\rotation$ agrees with the $2$-functor arising from the biadjunction of $[I \subset Ii]$ and $[Ii \supset I]$. After applying $\Phi$ one obtains the $2$-functor arising from the bijadunction of induction and restriction, with units and counits of adjunction determined by the Frobenius extension structures. In particular, this corollary implies that biadjunction on $\SBSBim$ agrees\footnote{Technically, this argument only implies agreement on the image of the functor $\Phi$. Recall that $\Phi$ was proven to be $2$-full for balanced realizations in \cite{EKLP}, and is only conjectured to be so for unbalanced realizations. The result does hold more generally though. Using the identifications $\vartheta_{IK_\bullet}:=\eta_{\Rev(IK_\bullet)} \circ \etadot_{IK_\bullet}:\BS(\Rev(IK_\bullet))\xrightarrow{\sim} \Dual \circ \Rev(\BS(IK_\bullet))$ 
%$\vartheta_{IK_\bullet}:\BS(\Rev(IK_\bullet))\xrightarrow{\sim} \Dual \circ \Rev(\BS(IK_\bullet))$ given by $\eta_{\Rev(IK_\bullet)} \circ \etadot_{IK_\bullet}$ 
for multistep expressions $IK_\bullet$, one can directly check that the resulting isomorphisms 
 \begin{gather*}
 \Hom_{\SBSBim}(\BS(IK_\bullet), \BS(IK'_\bullet)) \xrightarrow{\sim} \Hom_{\SBSBim}(\Rev (\BS(IK'_\bullet)), \Rev(\BS(IK_\bullet)))\\
 f \mapsto \vartheta_{IK_\bullet}^{-1}\circ \rotation(f) \circ \vartheta_{IK'_\bullet}
 \end{gather*}
  coincide precisely with the isomorphisms obtained by iteratively using the induction-restriction bi-adjunctions.} coherently with the composition of two other symmetries $\Dual$ and $\Rev$, which is not a priori obvious. 
\end{remark}


%\PT{Commented out the following unfinished remark.}
% \begin{remark} \BE{previous remark}
%     The isomorphisms induced by the 2-functor $\rotation$ on the 2-morphism spaces in $\SBSBim$ identifies with the isomorphisms one obtains from the induction-restriction bi-adjunctions. Note that these adjunction isomorphisms are uniquely determined only after fixing a choice of units and counits, and the latter is exactly obtained via the inclusion, multiplication, Frobenius trace and coproduct map for each $I \subset J$ of finitary subsets.
%     \PT{Add example.}
%     Showing that these isomorphisms on 2-morphisms can be combined coherently to make $\rotation$ into a 2-functor requires extra compatibility with 1-morphism composition, and this is exactly what's built into the definition of a Frobenius hypercube: for example, for $I\subset J \subset K$ of finitary subsets
%     \PT{Add}
%     is precisely because of the assumption that $\dd^J_K\circ\dd^I_J=\dd^I_K$ etc.
% \end{remark}

\begin{remark}\label{rmk left duals and right duals are naturally isomorphic}
    One may also define a \emph{left duality} 2-functor $\widetilde{\Dual}$ (see \cite[Remark 3.10]{BBE}) on $\SBSBim$ which fixes objects and act on 1-morphism categories via
    \begin{equation}
    \Hom(I,J)^{\op}\to \Hom(I,J): M \mapsto \Hom^\bullet_{R^I}(M, R^I)(2\ell(J)-2\ell(I))
    \end{equation}
    with suitable changes to the coherence 2-morphisms in \cref{eq coherence data for duality} and \cref{eq coherence data for duality}.
    This 2-functor is related to $\Dual$ by the equality (not just an isomorphism) of 2-functors
    \begin{equation}\label{eq left duality is conjugate to right duality by reverse}
        \tilde{D}=\Rev \circ \Dual \circ \Rev.
    \end{equation}
    Analogous to \cref{lemma self duality of Bott-Samelsons}, we have an isomorphism (canonical, after choosing Frobenius structures)
    \begin{gather}
        \notag \tilde{\eta}_{IK_\bullet}: \BS(IK_\bullet) \xrightarrow{\sim} \widetilde{\Dual}(\BS(IK_\bullet)) \\
        r_0\otimes \cdots r_d \mapsto \Bigl(s_0 \otimes \cdots \otimes s_d \mapsto s_0r_0 \;\dd^{I_1}_{K_1}\bigl(r_1s_1\cdots \dd^{I_{d-1}}_{K_{d-1}}(r_{d-1}s_{d-1}\dd^{I_d}_{K_d}(r_d s_d)) \cdots \bigr)\Bigr) \label{eq self duality of Bott Samelsons for left dual}
    \end{gather}
    for any singular multistep expression $IK_\bullet$.
    One can then check that the isomorphisms 
    \begin{equation}
    \eta_{IK_\bullet} \circ (\eta_{IK_\bullet})^{-1}: \Dual(\BS(IK_\bullet)) \rightarrow \widetilde{\Dual}(\BS(IK_\bullet))\end{equation}
    provide the data of an isomorphism of (weak) 2-functors $\Dual \Rightarrow \widetilde{\Dual}$.
   % Hence $\Phi$ also interwines $\Dual$ on $\mc{D}^{\pre}$ with $\widetilde{\Dual}$.
   Together with the equality in \cref{eq left duality is conjugate to right duality by reverse} and \cref{lemma dual and rev are involutions}, this also gives us an isomorphism of (weak) 2-functors $\Rev \circ \Dual \Rightarrow \Dual \circ \Rev$ (cf. \cref{corollary phi intertwines rotation}). 
\end{remark}

We now introduce certain non-degenerate bilinear forms on singular Bott-Samelson bimodules to compare with other approaches to duality in the literature.
Let $IK_\bullet$ be a singular multistep expression as in \cref{eq singular multistep expression supplementary}.
The self duality isomorphisms of singular Bott-Samelson bimodules (\cref{lemma self duality of Bott-Samelsons}) can be used to give a (right) $R^{I_d}$-bilinear form on $\BS(IK_\bullet)$ valued in $R^{I_d}$, using the canonical pairing
\begin{equation}
    \BS(IK_\bullet) \times \Dual(\BS(IK_\bullet)) \rightarrow R^{I_d}: (b, \beta)\mapsto \beta(b).
\end{equation}

\begin{defn}\label{defn invariant pairing on singular Bott Samelsons}
    Define the \emph{right (global) intersection form} on $\BS(IK_\bullet)$ to be the the (right) $R^{I_d}$-bilinear form given by 
    \begin{equation}
       \langle -,-\rangle_{IK_\bullet}: \BS(IK_\bullet)\times \BS(IK_\bullet) \rightarrow R^{I_d}: \langle b,b'\rangle_{IK_\bullet} := \bigl(\eta_{IK_\bullet}(b')\bigr)(b).
    \end{equation}
    From the explicit formula for $\eta_{IK\bullet}$ in \cref{eq self duality of bott samelsons}, 
    we have that when 
    $$b= r_0 \otimes \cdots \otimes r_d,\; b'=s_0\otimes \cdots \otimes s_d$$
    for $r_j, s_j \in R^{I_j}$, $0\leq j \leq d$,
    \begin{equation}
        \bigl(\eta_{IK_\bullet}(b')\bigr)(b) =\dd^{I_{d-1}}_{K_d} \left(\cdots \dd^{I_1}_{K_2}(\dd^{I_0}_{K_1}(r_0s_0)r_1s_1) \cdots r_{d-1}s_{d-1} \right)r_ds_d= \bigl(\eta_{IK_\bullet}(b)\bigr)(b')
    \end{equation}
    so that $\langle-,-\rangle_{IK_\bullet}$ is a symmetric bilinear form. 
    Moreover, it is $(R^{I_0}, R^{I_d})$-invariant, i.e., for any $f\in R^{I_0}$ and $g \in R^{I_d}$, we have
    \begin{equation}
        \langle fmg,m' \rangle_{IK_\bullet}=\langle m, fm'g\rangle_{IK_\bullet}.
    \end{equation}
    This form is also non-degenerate, since $\eta_{IK_\bullet}$ is an isomorphism and $\BS(IK_\bullet)$ is a free right $R^{I_d}$-module.
\end{defn}

\begin{remark}\label{rmk left intersection form}
    One can similarly also define a \emph{left (global) intersection form} $\widetilde{\langle -, -\rangle}_{IK_\bullet}$ on $\BS(IK_\bullet)$, which will be a symmetric $(R^{I_0},R^{I_d})$-invariant non-degenerate (left) $R^{I_0}$-bilinear form valued in $R^{I_0}$, using the self duality isomorphisms \cref{eq self duality of Bott Samelsons for left dual} for the left dual $\widetilde{D}$ in \cref{rmk left duals and right duals are naturally isomorphic}. 
    This form is related to the right intersection form by the equation
    \begin{equation}
        \widetilde{\langle f, g \rangle}_{IK_\bullet}= \langle \etadot_{IK_\bullet}^{-1}(f), \etadot_{IK_\bullet}^{-1}(g)\rangle_{\Rev (IK_\bullet)}
    \end{equation}
    where $\Rev(IK_\bullet)$ is the reverse of $IK_\bullet$ as in \cref{defn reverse of singular multistep expression} and $\etadot$ is as in \cref{lemma reversal on bott samelsons and multistep expressions}.
\end{remark}

\begin{remark}\label{rmk intersection forms in literature}
In Williamson's thesis \cite{WillThesis} where singular Soergel bimodules are originally defined, it is shown that the functor of tensoring with an induction or restriction bimodule commutes with duality, and the natural isomorphism involved is defined using the Frobenius structure, precisely as  $\eta_{I_\bullet}$ for the cases when $I_\bullet$ equals $[I \subset Ii]$ or $[Ii \supset I]$.  This happens in \cite[\S 3.2]{WillThesis} and \cite[Proposition 4.3.4]{WillThesis}, though the details are omitted from the published article \cite{WillSingular}. Iterating this natural isomorphism over iterated tensor products, one obtains $\eta_{I_\bullet}$ for any singlestep expression $I_\bullet$. This approach is also taken in \cite{BBE}.

Meanwhile, the subsequent literature tended to define the isomorphism between a Bott-Samelson bimodule and its dual using the global intersection form. References include \cite{EWHodge} and \cite{EMTW} for ordinary Soergel bimodules, and \cite[Appendix A]{patimo2021basesintersectioncohomologygrassmannian} for one-sided-singular Soergel bimodules, though we are unaware of a reference which treats global intersection forms for singular Soergel bimodules in general. These papers define the intersection form inductively, having chosen a particular top-degree element of each induction bimodule. In \cite[Appendix A]{patimo2021basesintersectioncohomologygrassmannian}, the Frobenius structure is explicitly used to choose a top-degree element.


% Thus the relationship between duality and the Frobenius structure is present but somewhat obfuscated by the literature, which also obfuscates why there should be a relationship between rotation and flipping. 


    

    
    % While there are many references in the literature for the global intersection form on ordinary (non-singular) Bott-Samelson bimodules (\cite{EWHodge}, \cite{EMTW}), as well as extensions to one-sided singular Bott Samelson bimodules (\cite[Appendix A]{patimo2021basesintersectioncohomologygrassmannian}), the authors are not aware of a proper reference for the intersection forms on singular Bott-Samelson bimodules in full generality. \BE{edits}
    % Most \PT{all?} of the existing constructions involve an inductive argument to construct this pairing, while we provide a completely explicit formula. These do agree, for example, by \cite[Proposition A.12]{patimo2021basesintersectioncohomologygrassmannian} and \cref{thm phi commutes with dual and rev}.

\end{remark}

\begin{remark}
    While we explicitly defined the 2-functors $\Dual$ and $\Rev$ only on $\SBSBim$, it is straightforward to extend these 2-functors, as well as the results involving them, to the Karoubi envelope $\SSBim$.
\end{remark}


\subsection{Light ladders and light leaves: an example} \label{subsec:LLappendix}

As we illustrate a particular example, we briefly recall the purpose of elementary light leaves and ladders.

We consider representations of $\slf_n$, and let $V_k := \Lambda^k(\C^n)$. Let $\mu$ be a weight in $V_k$ for some $k$. There is a minimal dominant weight $\lambda_{\start}$ such that $\lambda_{\finish} := \lambda_{\start} + \mu$ is also dominant. If $L_{\lambda}$ denotes the simple representation of highest weight $\lambda$, then $L_{\lambda_{\finish}}$ is a direct summand of $L_{\lambda_{\start}} \ot V_k$, and we are interested in constructing the projection map.

Now $L_{\lambda_{\start}}$ is a summand of a tensor product $V_{i_1} \ot \cdots \ot V_{i_d}$ of fundamental representations, when $\sum_{j=1}^d \varpi_{i_j} = \lambda_{\start}$. Similarly, $L_{\lambda_{\finish}}$ is a summand of $V_{j_1} \ot \cdots \ot V_{j_e}$. The \emph{elementary light ladder} $\ELL(\mu)$ is a particular morphism
\[ \ELL(\mu) \colon (V_{i_1} \ot \cdots \ot V_{i_d}) \ot V_k \to (V_{j_1} \ot \cdots \ot V_{j_e}) \]
which remains nonzero after pre- and post-composition with the appropriate idempotents, so that it descends to a nonzero (projection) map
\[ L_{\lambda_{\start}} \ot V_k \to L_{\lambda_{\finish}}.\]

\begin{remark} \label{rmk:whenlambdanotminimal} If $\lambda$ is any dominant weight such that $\lambda + \mu$ is dominant, then $\lambda - \lambda_{\start}$ is also dominant. The projection map $L_{\lambda} \ot V_k \to L_{\lambda + \mu}$ can be obtained from $\id \ot \ELL(\mu)$ by pre- and post-composition with the appropriate idempotents. Here $\id$ represents the identity map of a tensor product of fundamental representations associated with $\lambda - \lambda_{\start}$. \end{remark}

Here is a particular example when $n \ge 5$ and $k=2$. Let $\mu = \varpi_4 - \varpi_3 + \varpi_2 - \varpi_1$. Then $\lambda_{\start} = \varpi_1 + \varpi_3$, and $\lambda_{\finish} = \varpi_2 + \varpi_4$. Note that $L_{\lambda_{\start}}$ is a summand of $V_1 \ot V_3$ and $L_{\lambda_{\finish}}$ is a summand of $V_2 \ot V_4$. Then 
\begin{equation} \ELL(\mu) \colon (V_1 \ot V_3) \ot V_2 \to (V_2 \ot V_4), \qquad \ELL(\mu) = \isvg{0.15}{ELLforBen2}. \end{equation}
We have chosen, arbitrarily, to label the rightmost region with the label $0 \in \Omega$. After applying $\GS$ or $\GS_{\zeta}$ we obtain
\begin{equation} \label{eq:GSELLdiagram} \GS(\ELL(\mu)) = \textrm{scalar} \cdot {
\labellist
\small\hair 2pt
 \pinlabel {$\blk{5}$} [ ] at 8 -5
 \pinlabel {$\blk{6}$} [ ] at 24 -5
 \pinlabel {$\blk{2}$} [ ] at 40 -5
 \pinlabel {$\blk{5}$} [ ] at 56 -5
 \pinlabel {$\blk{0}$} [ ] at 72 -5
 \pinlabel {$\blk{2}$} [ ] at 88 -5
 \pinlabel {$\blk{4}$} [ ] at 8 80
 \pinlabel {$\blk{6}$} [ ] at 24 80
 \pinlabel {$\blk{0}$} [ ] at 72 80
 \pinlabel {$\blk{4}$} [ ] at 88 80
 \pinlabel {$\dgr{0}$} [ ] at 90 45
\endlabellist
\centering
\ig{1}{NastyELLimage}
}. \end{equation}

\vspace{.1cm}

We note that the colors $1, 3 \in \Omega$ remain present in every region of this diagram.

Now we consider elementary light leaves for singular Soergel bimodules. Whenever $I$ and $J \cup s$ are finitary, for $s \notin J$, and whenever $p \in W_I \backslash W / W_J$ and $q \in W_I \backslash W / W_{Js}$ are such that $p \subset q$, we call $[p \subset q]$ or $[q \supset p]$ a \emph{coset pair}. Observe that $B_q$ is a summand of $B_p \ot \BS([J \subset Js])$, and $B_p$ is a summand of $B_q \ot \BS([Js \supset J])$. We are interested in constructing the projection maps.

Choosing a reduced expressions $IK_{\bullet}$ for $p$, we have that $B_p$ is a summand of $\BS(IK_{\bullet})$. Similarly for a reduced expression $IK'_{\bullet}$ of $q$. The elementary light ladder $\ELL([p \subset q])$ is a particular morphism
\[ \ELL([p \subset q]) \colon \BS(IK_{\bullet}) \ot \BS(J \subset Js) \to \BS(IK'_{\bullet})\]
which descends to a nonzero projection map $B_p \ot \BS(J \subset Js) \to B_q$. Similarly, 
\[ \ELL([q \supset p]) \colon \BS(IK'_{\bullet}) \ot \BS(Js \supset J) \to \BS(IK_{\bullet}) \] descends to a projection map $B_q \ot \BS(Js \supset J) \to B_p$.

Soon enough we will wish to illustrate double cosets, for which it helps to fix a particular value of $n$, and we choose $n=8$. Below we will be using the symbol $n$ to indicate a particular double coset rather than the size of $\Omega$, following the notation in \cite{EKLP}. 

Just as in Remark \ref{rmk:whenlambdanotminimal}, $[p \subset q]$ might not be ``minimal.'' In this case there is a coset $z$ (think of $z$ as the analogue of $\lambda - \lambda_{\start}$) and a \emph{Grassmannian coset pair} $[m \subset n]$ such that $z.m = p$ and $z.n = q$ (i.e. a reduced expression for $z$ concatenated with a reduced expression for $m$ is a reduced expression for $p$). One has
\begin{equation} \ELL([p \subset q]) := \id_z \ot \ELL([m \subset n]), \end{equation}
where by $\id_z$ we mean the identity map of a reduced expression for $z$.

The analogue of $(-)\ot V_k$ after applying $\GS$ will be $(-) \ot [\hh{k} \supset \hh{0,k} \subset \hh{0}] = (-) \ot [-0+k]$. This involves two steps, not one, so we are interested in a composition of two elementary light leaves. Starting with a coset $q \in W_I \backslash W / W_{\hh{k}}$ we would find $[q \supset p \subset q']$ where
\[ p \in W_I \backslash W / W_{\hh{0,k}}, \qquad q' \in W_I \backslash W / W_{\hh{0}}\]
and take the composition $\ELL([p \subset q']) \circ (\ELL([q \supset p]) \ot \id_{[+k]})$.

Let 
\[ q_{\finish} = \psi_0(\lambda_{\finish}) \in W_{\hh{6}} \backslash W / W_{\hh{0}}, \qquad \text{resp. } \; q_{\start} = \psi_2(\lambda_{\start}) \in W_{\hh{6}} \backslash W / W_{\hh{2}}.\]
Following the bijection of \cref{def:psia} we have
\begin{equation} \overline{q}_{\finish} w_{\hh{0}} = {
\labellist
\tiny\hair 2pt
 \pinlabel {$1$} [ ] at 120 0
 \pinlabel {$6$} [ ] at 160 45
\endlabellist
\centering
\ig{1}{qendw0}
}, \end{equation}
\begin{equation} \overline{q}_{\start} w_{\hh{2}} = {
\labellist
\tiny\hair 2pt
 \pinlabel {$1$} [ ] at 120 0
 \pinlabel {$6$} [ ] at 160 45
\endlabellist
\centering
\ig{1}{qstartw0}
}. \end{equation}
Let us also provide the minimal coset representatives:
\begin{equation} \label{eq:qendmin} \underline{q}_{\finish} = {
\labellist
\tiny\hair 2pt
 \pinlabel {$1$} [ ] at 120 0
 \pinlabel {$6$} [ ] at 160 45
\endlabellist
\centering
\ig{1}{qendmin}
}, \end{equation}
\begin{equation} \underline{q}_{\start} = {
\labellist
\tiny\hair 2pt
 \pinlabel {$1$} [ ] at 120 0
 \pinlabel {$6$} [ ] at 160 45
\endlabellist
\centering
\ig{1}{qstartmin}
}. \end{equation}
(Recall that $\overline{q}$ is the maximal coset representative and $\underline{q}$ is the minimal coset representative, in the Bruhat order.)

There is a unique double coset $p \in W_{\hh{6}} \backslash W / W_{\hh{0,2}}$ for which $p \subset q_{\start} \cap q_{\finish}$, and we have
\begin{equation} \underline{p} = {
\labellist
\tiny\hair 2pt
 \pinlabel {$1$} [ ] at 120 0
 \pinlabel {$6$} [ ] at 160 45
\endlabellist
\centering
\ig{1}{pmin}
}.\end{equation}

We wish to construct $\ELL([p \subset q_{\finish}])$, following the algorithm in \cite[\S 7]{EKLP}.

The right redundancy of $q_{\finish}$ is $Q = Q_{\finish} = \hh{024}$. Then $z = z_{\finish} \in W_{\hh{6}} \backslash W / W_Q$ satisfies $\underline{z} = \underline{q}_{\finish}$. A reduced expression for $z$ (in two different notations) is
\begin{equation} z \expr [[\hh{6} \supset \hh{246} \subset \hh{24} \supset \hh{024}]] = [\hh{6} - 2 - 4 + 6 - 0]. \end{equation}
One should think that $[\hh{6} - 2 - 4]$ breaks the parabolic subgroup into blocks associated to the different colors in \eqref{eq:qendmin}, and $[+6 - 0]$ is the cabled crossing of the red and blue strands.

Continuing, we have $m = m_{\finish} \in W_Q \backslash W / W_{\hh{02}}$ with $\underline{q}_{\finish} \cdot \underline{m} = \underline{p}$. Here is the picture:
\begin{equation} \underline{m} = {
\labellist
\tiny\hair 2pt
 \pinlabel {$1$} [ ] at 120 0
 \pinlabel {$6$} [ ] at 160 45
\endlabellist
\centering
\ig{1}{mend}
}. \end{equation}
The reader should be able to see that stacking $\underline{q}_{\finish}$ above $\underline{m}_{\finish}$ yields $\underline{p}$. A reduced expression for $m$ is
\begin{equation} m \expr [Q - 1 - 3 + 2 - 2 + 1 + 3 + 4]. \end{equation}
The portion $[+2 - 2]$ of this expression is responsible for the crossing.
Now $n = n_{\finish} \in W_Q \backslash W / W_{\hh{0}}$ has reduced expression
\begin{equation} n \expr [Q + 2 + 4]. \end{equation}
The pair $[m \subset n]$ is a Grassmannian coset pair for $I = \hh{6}$ and $J = \hh{02}$ and $s = s_2$.

Tensoring $\ELL([m \subset n])$ (a ``sinister'' diagram with lots of left crossings and leftwards caps) with $\id_z$ yields
\begin{equation} \ELL([p \subset q_{\finish}]) = {
\labellist
\small\hair 2pt
 \pinlabel {$\blk{4}$} [ ] at 8 -4
 \pinlabel {$\blk{2}$} [ ] at 24 -4
 \pinlabel {$\blk{6}$} [ ] at 40 -4
 \pinlabel {$\blk{0}$} [ ] at 56 -4
 \pinlabel {$\blk{1}$} [ ] at 72 -4
 \pinlabel {$\blk{3}$} [ ] at 88 -4
 \pinlabel {$\blk{2}$} [ ] at 104 -4
 \pinlabel {$\blk{2}$} [ ] at 120 -4
 \pinlabel {$\blk{3}$} [ ] at 136 -4
 \pinlabel {$\blk{1}$} [ ] at 152 -4
 \pinlabel {$\blk{4}$} [ ] at 168 -4
 \pinlabel {$\blk{2}$} [ ] at 184 -4
 \pinlabel {$\blk{4}$} [ ] at 8 71
 \pinlabel {$\blk{2}$} [ ] at 24 71
 \pinlabel {$\blk{6}$} [ ] at 40 71
 \pinlabel {$\blk{0}$} [ ] at 56 71
 \pinlabel {$\blk{2}$} [ ] at 72 71
 \pinlabel {$\blk{4}$} [ ] at 120 71
 \pinlabel {$\dgr{0}$} [ ] at 184 45
\endlabellist
\centering
\ig{1}{ELLend}
}. \end{equation}

\vspace{.1cm}

Forgive us as we reuse the notation $Q, z, m, n$ for a new case. The right redundancy of $q_{\start}$ is $Q = Q_{\start} = \hh{235}$. Then $z = z_{\start} \in W_{\hh{6}} \backslash W / W_Q$ satisfies $\underline{z} = \underline{q}_{\start}$. A reduced expression for $z$ is
\begin{equation} z \expr [[\hh{6} \supset \hh{3,5,6} \subset \hh{35} \supset \hh{235}]] = [\hh{6} - 5 - 3 + 6 - 2]. \end{equation}
Again, $[\hh{6} - 5 - 3]$ breaks the parabolic subgroup into blocks associated to the different colors in the picture for $\underline{q}_{\start}$, and $[+6 - 2]$ is the cabled crossing of the red and blue strands.

Continuing, we have $m = m_{\start} \in W_Q \backslash W / W_{\hh{02}}$ with $\underline{q}_{\start} \cdot \underline{m} = \underline{p}$. Here is the picture:
\begin{equation} \underline{m} = {
\labellist
\tiny\hair 2pt
 \pinlabel {$1$} [ ] at 120 0
 \pinlabel {$6$} [ ] at 160 45
\endlabellist
\centering
\ig{1}{mstart}
}. \end{equation}
The reader should be able to see that stacking $\underline{q}_{\start}$ above $\underline{m}_{\start}$ yields $\underline{p}$. A reduced expression for $m$ is
\begin{equation} m \expr [[Q - 1 - 4 + 5 - 0 + 1 + 3 + 4]]. \end{equation}
The portion $[+5 - 0]$ of this expression is responsible for the cabled crossing.
Now $n = n_{\start} \in W_Q \backslash W / W_{\hh{2}}$ has reduced expression
\begin{equation} n \expr [Q + 3 + 5]. \end{equation}
The pair $[m \subset n]$ is a Grassmannian coset pair for $I = \hh{6}$ and $J = \hh{02}$ and $s = s_0$.

One obtains $\ELL([n \supset m])$ by flipping $\ELL([m \subset n])$ upside-down and using adjunction. Tensoring $\ELL([n \supset m])$ with $\id_z$ yields
\begin{equation} \ELL([q_{\start} \supset p]) = {
\labellist
\small\hair 2pt
 \pinlabel {$\blk{5}$} [ ] at 8 -2
 \pinlabel {$\blk{3}$} [ ] at 24 -2
 \pinlabel {$\blk{6}$} [ ] at 40 -2
 \pinlabel {$\blk{2}$} [ ] at 56 -2
 \pinlabel {$\blk{5}$} [ ] at 72 -2
 \pinlabel {$\blk{3}$} [ ] at 88 -2
 \pinlabel {$\blk{0}$} [ ] at 168 -2
 \pinlabel {$\blk{5}$} [ ] at 8 71
 \pinlabel {$\blk{3}$} [ ] at 24 71
 \pinlabel {$\blk{6}$} [ ] at 40 71
 \pinlabel {$\blk{2}$} [ ] at 56 71
 \pinlabel {$\blk{1}$} [ ] at 72 71
 \pinlabel {$\blk{4}$} [ ] at 88 71
 \pinlabel {$\blk{5}$} [ ] at 104 71
 \pinlabel {$\blk{0}$} [ ] at 120 71
 \pinlabel {$\blk{3}$} [ ] at 136 71
 \pinlabel {$\blk{4}$} [ ] at 152 71
 \pinlabel {$\blk{1}$} [ ] at 168 71
\pinlabel {$\dgr{02}$} [ ] at 168 35
\endlabellist
\centering
\ig{1}{ELLstart}
}. \end{equation}

Now we wish to compose $\ELL([q_{\finish} \supset p])$ with $\ELL([p \subset q_{\start}]) \otimes \id_{[+2]}$ to obtain a morphism comparable with $\GS(\ELL(\mu))$ from \eqref{eq:GSELLdiagram}. However, the sources and targets of these morphisms do not match, because they use different reduced expressions for $q_{\start}$ and $q_{\finish}$ and $p$. These different reduced expressions were essential aspects of the algorithm, since at various stages we were required to use reduced expressions starting with $z_{\start}$ or $z_{\finish}$. Different reduced expressions for the same double coset are related by the (singular) braid relations, introduced in \cite{EKo}, and to each braid relation we have a corresponding diagrammatic morphism called a rex move, introduced in \cite[\S 6]{EKLP}.
To compose these elementary light leaves we should insert rex moves in between.

Examples of rex moves are: crossings of upward-oriented strands (the up-up move), crossings of downward oriented strands (down-down), sideways crossings where the colors commute in the ambient parabolic subgroup (commuting switchback), and the type $A$ switchback maps which look like a cup or cap straddling two other strands, see \cite[Example 6.1]{EKLP}. See also \cref{rmk:switchback}.


A rex move from the target of $\ELL([p \subset q_{\finish}])$ to the target of $\GS(\ELL(\mu))$ is a single type $A$ switchback
\begin{equation} \rex_3 = \ig{1}{rex3}. \end{equation}
As a reminder:
\[ \ELL([p \subset q_{\finish}]) = \ig{1}{ELLend}. \]
A rex move from the target of $\ELL([q_{\start} \supset p]) \ot \id_{[+2]}$ to the source of $\ELL([p \subset q_{\finish}])$ is 
\begin{equation} \rex_2 = \ig{1}{rex2}. \end{equation}
This is a composition of all the rex moves mentioned elsewhere.
As a reminder:
\[ \ELL([q_{\start} \supset p]) \ot \id_{[+2]} = 
\ig{1}{ELLstartwid}
. \]
A rex move from the source of $\GS(\ELL(\mu))$ to the source of $\ELL([q_{\start} \supset p]) \ot \id_{[+2]}$ is
\begin{equation} \rex_1 = \ig{1}{rex1}. \end{equation}
The orange cup straddling the red and fuschia strands is a type $A$ switchback, and the orange-green crossing is up-up.


\begin{remark} \label{rmk:switchback} Let $\Dual$ flip diagrams upside-down and reverse orientations without changing region labels. The composition $\rex_3 \circ \Dual(\rex_3)$ is equal to the identity map. This is called the \emph{switchback relation} and is proven in \cite[Theorem 6.7]{EKLP} as a consequence of the relations from \cite{EWSFrob}. A prototypical depiction of this relation is \cite[Lemma 3.4]{BBE}. \end{remark}

Composing all these morphisms together, we obtain
\begin{equation} \label{eq:finalnastydiagram} \rex_3 \circ \ELL([p \subset q_{\finish}]) \circ \rex_2 \circ (\ELL([q_{\start} \supset p]) \ot \id_{[+2]}) \circ \rex_1 = \ig{1.2}{NastyLL}. \end{equation}
From here, one can use a series of the relations: oriented R3 \eqref{R3oriented}, distant R3 \eqref{distantR3}, distant R2 \eqref{R2 distant}, and switchback (see \cref{rmk:switchback}), eventually deducing that the diagrams in \eqref{eq:finalnastydiagram} and \eqref{eq:GSELLdiagram} are equal. We leave this exercise to the reader.

We note that the colors $1, 3 \in \Omega$ (navy blue and orange) do not label strands in \eqref{eq:GSELLdiagram}. They arise in \eqref{eq:finalnastydiagram} because the light leaf must factor through a reduced expression for $p$.
